% paper64-team-workflow.tex
% Multi-Developer Workflows with Jurisdiction and Authority Delegation
% JuGeo Paper Series — Paper 64

\documentclass[11pt]{article}
% jugeo-common.tex — shared preamble for all Judgment Geometry papers
% Include via: % jugeo-common.tex — shared preamble for all Judgment Geometry papers
% Include via: % jugeo-common.tex — shared preamble for all Judgment Geometry papers
% Include via: \input{jugeo-common}

% ── Packages ────────────────────────────────────────────────────────────
\usepackage{amsmath,amssymb,amsthm,mathtools}
\usepackage{tikz-cd}
\usepackage{listings}
\usepackage{xcolor}
\usepackage{xspace}
\usepackage{booktabs}
\usepackage{multirow}
\usepackage{enumitem}
\usepackage[expansion=false]{microtype}
\usepackage{hyperref}
\usepackage{cleveref}
% \usepackage{thmtools}  % disabled: not available in this TeX installation
\usepackage{float}
\usepackage{caption}
\usepackage{subcaption}
\usepackage{tabularx}
\usepackage{stmaryrd}       % \llbracket, \rrbracket
\usepackage{bbm}            % \mathbbm{1}

% ── Page geometry ───────────────────────────────────────────────────────
\usepackage[margin=1in]{geometry}

% ── Theorem environments ────────────────────────────────────────────────
\theoremstyle{plain}
\newtheorem{theorem}{Theorem}[section]
\newtheorem{lemma}[theorem]{Lemma}
\newtheorem{proposition}[theorem]{Proposition}
\newtheorem{corollary}[theorem]{Corollary}

\theoremstyle{definition}
\newtheorem{definition}[theorem]{Definition}
\newtheorem{example}[theorem]{Example}
\newtheorem{construction}[theorem]{Construction}

\theoremstyle{remark}
\newtheorem{remark}[theorem]{Remark}
\newtheorem{notation}[theorem]{Notation}

% ── Python listings style ──────────────────────────────────────────────
\definecolor{jugeoBlue}{HTML}{2563EB}
\definecolor{jugeoGreen}{HTML}{059669}
\definecolor{jugeoRed}{HTML}{DC2626}
\definecolor{jugeoPurple}{HTML}{7C3AED}
\definecolor{jugeoGray}{HTML}{6B7280}
\definecolor{jugeoBg}{HTML}{F8FAFC}

\lstdefinestyle{jugeo-python}{
  language=Python,
  basicstyle=\ttfamily\small,
  keywordstyle=\color{jugeoBlue}\bfseries,
  stringstyle=\color{jugeoGreen},
  commentstyle=\color{jugeoGray}\itshape,
  numberstyle=\tiny\color{jugeoGray},
  backgroundcolor=\color{jugeoBg},
  numbers=left,
  numbersep=6pt,
  frame=single,
  framerule=0.4pt,
  rulecolor=\color{jugeoGray!40},
  breaklines=true,
  breakatwhitespace=true,
  showstringspaces=false,
  tabsize=4,
  xleftmargin=1.5em,
  framexleftmargin=1.5em,
  morekeywords={dataclass,frozen,field,Enum,IntEnum,auto,Any,
                Mapping,Sequence,Optional,match,case},
  literate={→}{{$\to$}}1 {←}{{$\leftarrow$}}1
           {≤}{{$\leq$}}1 {≥}{{$\geq$}}1 {≠}{{$\neq$}}1
           {∈}{{$\in$}}1 {∀}{{$\forall$}}1 {∃}{{$\exists$}}1
           {⊕}{{$\oplus$}}1 {⊖}{{$\ominus$}}1,
  escapeinside={(*@}{@*)},
}
\lstset{style=jugeo-python}

\lstdefinestyle{jugeo-output}{
  basicstyle=\ttfamily\small,
  backgroundcolor=\color{jugeoBg},
  frame=single,
  framerule=0.4pt,
  rulecolor=\color{jugeoGray!40},
  breaklines=true,
  showstringspaces=false,
  xleftmargin=1.5em,
  framexleftmargin=0.5em,
  numbers=none,
}

% ── JuGeo-specific macros ──────────────────────────────────────────────

% Judgment 8-tuple
\newcommand{\J}{\mathcal{J}}
\newcommand{\Jud}[8]{(#1,\,#2,\,#3,\,#4,\,#5,\,#6,\,#7,\,#8)}
\newcommand{\JudShort}[1]{\J_{#1}}

% Components
\newcommand{\coord}{\mathsf{c}}            % coordinate
\newcommand{\prop}{\varphi}                 % proposition
\newcommand{\carrier}{\mathcal{A}}          % carrier type
\newcommand{\evid}{\mathcal{E}}             % evidence bundle
\newcommand{\oblig}{\mathcal{O}}            % obligations
\newcommand{\obstr}{\mathcal{B}}            % obstructions
\newcommand{\trust}{\mathcal{T}}            % trust annotation
\newcommand{\prov}{\Pi}                     % provenance

% Trust algebra
\newcommand{\Talg}{\mathfrak{T}}            % trust ordered algebra
\newcommand{\Eadm}{\mathcal{E}_{\mathrm{adm}}}
\newcommand{\tleq}{\preceq}                 % trust ordering
\newcommand{\tjoin}{\oplus}                 % conservative join
\newcommand{\tatten}{\ominus}               % attenuation
\newcommand{\tpromote}{\uparrow_\pi}        % promotion
\newcommand{\tdemote}{\downarrow_\chi}       % demotion

% Trust tiers
\newcommand{\Tcontra}{\ensuremath{\mathsf{CONTRADICTED}}}
\newcommand{\Tunver}{\ensuremath{\mathsf{UNVERIFIED}}}
\newcommand{\Tcopilot}{\ensuremath{\mathsf{COPILOT}}}
\newcommand{\Toracle}{\ensuremath{\mathsf{ORACLE}}}
\newcommand{\Truntime}{\ensuremath{\mathsf{RUNTIME}}}
\newcommand{\Tsolver}{\ensuremath{\mathsf{SOLVER}}}
\newcommand{\Tproof}{\ensuremath{\mathsf{PROOF}}}

% Site and geometry
\newcommand{\Site}{\mathcal{S}}             % semantic site
\newcommand{\Cov}{\mathrm{Cov}}            % covering family
\newcommand{\Ob}{\mathrm{Ob}}              % objects
\newcommand{\Mor}{\mathrm{Mor}}            % morphisms
\newcommand{\res}{\mathrm{res}}            % restriction
\newcommand{\Cech}{\check{C}}              % Čech complex
\newcommand{\Hcoh}[1]{H^{#1}}             % cohomology group

% Descent
\newcommand{\descent}{\mathrm{desc}}
\newcommand{\glue}{\mathrm{glue}}
\newcommand{\overlap}[2]{#1 \cap #2}

% Semantic moves
\newcommand{\VERIFY}{\textsc{Verify}}
\newcommand{\CONSTRUCT}{\textsc{Construct}}
\newcommand{\NORMALIZE}{\textsc{Normalize}}
\newcommand{\GLUE}{\textsc{Glue}}
\newcommand{\RESTRICT}{\textsc{Restrict}}
\newcommand{\TRANSPORT}{\textsc{Transport}}
\newcommand{\REPAIR}{\textsc{Repair}}
\newcommand{\PROMOTE}{\textsc{Promote}}
\newcommand{\CHALLENGE}{\textsc{Challenge}}
\newcommand{\ESCALATE}{\textsc{Escalate}}

% Fragments
\newcommand{\QFLIA}{\texttt{QF\_LIA}}
\newcommand{\QFLRA}{\texttt{QF\_LRA}}
\newcommand{\QFBV}{\texttt{QF\_BV}}
\newcommand{\QFUF}{\texttt{QF\_UF}}

% Misc
\newcommand{\Python}{\textsc{Python}}
\newcommand{\JuGeo}{\textsc{JuGeo}}
\newcommand{\LEAN}{\textsc{Lean}\,4}
\newcommand{\Fstar}{F$^{\star}$}
\newcommand{\Coq}{\textsc{Coq}}
\newcommand{\Dafny}{\textsc{Dafny}}

% Common shortcuts
\newcommand{\ie}{i.e.\@\xspace}
\newcommand{\eg}{e.g.\@\xspace}
\newcommand{\cf}{cf.\@\xspace}
\newcommand{\wrt}{w.r.t.\@\xspace}

% Evaluation shortcuts
\newcommand{\numcases}{300}
\newcommand{\numspec}{100}
\newcommand{\numeq}{100}
\newcommand{\numbug}{100}
\newcommand{\medtime}{6.5\,ms}
\newcommand{\totaltime}{1.949\,s}


% ── Packages ────────────────────────────────────────────────────────────
\usepackage{amsmath,amssymb,amsthm,mathtools}
\usepackage{tikz-cd}
\usepackage{listings}
\usepackage{xcolor}
\usepackage{xspace}
\usepackage{booktabs}
\usepackage{multirow}
\usepackage{enumitem}
\usepackage[expansion=false]{microtype}
\usepackage{hyperref}
\usepackage{cleveref}
% \usepackage{thmtools}  % disabled: not available in this TeX installation
\usepackage{float}
\usepackage{caption}
\usepackage{subcaption}
\usepackage{tabularx}
\usepackage{stmaryrd}       % \llbracket, \rrbracket
\usepackage{bbm}            % \mathbbm{1}

% ── Page geometry ───────────────────────────────────────────────────────
\usepackage[margin=1in]{geometry}

% ── Theorem environments ────────────────────────────────────────────────
\theoremstyle{plain}
\newtheorem{theorem}{Theorem}[section]
\newtheorem{lemma}[theorem]{Lemma}
\newtheorem{proposition}[theorem]{Proposition}
\newtheorem{corollary}[theorem]{Corollary}

\theoremstyle{definition}
\newtheorem{definition}[theorem]{Definition}
\newtheorem{example}[theorem]{Example}
\newtheorem{construction}[theorem]{Construction}

\theoremstyle{remark}
\newtheorem{remark}[theorem]{Remark}
\newtheorem{notation}[theorem]{Notation}

% ── Python listings style ──────────────────────────────────────────────
\definecolor{jugeoBlue}{HTML}{2563EB}
\definecolor{jugeoGreen}{HTML}{059669}
\definecolor{jugeoRed}{HTML}{DC2626}
\definecolor{jugeoPurple}{HTML}{7C3AED}
\definecolor{jugeoGray}{HTML}{6B7280}
\definecolor{jugeoBg}{HTML}{F8FAFC}

\lstdefinestyle{jugeo-python}{
  language=Python,
  basicstyle=\ttfamily\small,
  keywordstyle=\color{jugeoBlue}\bfseries,
  stringstyle=\color{jugeoGreen},
  commentstyle=\color{jugeoGray}\itshape,
  numberstyle=\tiny\color{jugeoGray},
  backgroundcolor=\color{jugeoBg},
  numbers=left,
  numbersep=6pt,
  frame=single,
  framerule=0.4pt,
  rulecolor=\color{jugeoGray!40},
  breaklines=true,
  breakatwhitespace=true,
  showstringspaces=false,
  tabsize=4,
  xleftmargin=1.5em,
  framexleftmargin=1.5em,
  morekeywords={dataclass,frozen,field,Enum,IntEnum,auto,Any,
                Mapping,Sequence,Optional,match,case},
  literate={→}{{$\to$}}1 {←}{{$\leftarrow$}}1
           {≤}{{$\leq$}}1 {≥}{{$\geq$}}1 {≠}{{$\neq$}}1
           {∈}{{$\in$}}1 {∀}{{$\forall$}}1 {∃}{{$\exists$}}1
           {⊕}{{$\oplus$}}1 {⊖}{{$\ominus$}}1,
  escapeinside={(*@}{@*)},
}
\lstset{style=jugeo-python}

\lstdefinestyle{jugeo-output}{
  basicstyle=\ttfamily\small,
  backgroundcolor=\color{jugeoBg},
  frame=single,
  framerule=0.4pt,
  rulecolor=\color{jugeoGray!40},
  breaklines=true,
  showstringspaces=false,
  xleftmargin=1.5em,
  framexleftmargin=0.5em,
  numbers=none,
}

% ── JuGeo-specific macros ──────────────────────────────────────────────

% Judgment 8-tuple
\newcommand{\J}{\mathcal{J}}
\newcommand{\Jud}[8]{(#1,\,#2,\,#3,\,#4,\,#5,\,#6,\,#7,\,#8)}
\newcommand{\JudShort}[1]{\J_{#1}}

% Components
\newcommand{\coord}{\mathsf{c}}            % coordinate
\newcommand{\prop}{\varphi}                 % proposition
\newcommand{\carrier}{\mathcal{A}}          % carrier type
\newcommand{\evid}{\mathcal{E}}             % evidence bundle
\newcommand{\oblig}{\mathcal{O}}            % obligations
\newcommand{\obstr}{\mathcal{B}}            % obstructions
\newcommand{\trust}{\mathcal{T}}            % trust annotation
\newcommand{\prov}{\Pi}                     % provenance

% Trust algebra
\newcommand{\Talg}{\mathfrak{T}}            % trust ordered algebra
\newcommand{\Eadm}{\mathcal{E}_{\mathrm{adm}}}
\newcommand{\tleq}{\preceq}                 % trust ordering
\newcommand{\tjoin}{\oplus}                 % conservative join
\newcommand{\tatten}{\ominus}               % attenuation
\newcommand{\tpromote}{\uparrow_\pi}        % promotion
\newcommand{\tdemote}{\downarrow_\chi}       % demotion

% Trust tiers
\newcommand{\Tcontra}{\ensuremath{\mathsf{CONTRADICTED}}}
\newcommand{\Tunver}{\ensuremath{\mathsf{UNVERIFIED}}}
\newcommand{\Tcopilot}{\ensuremath{\mathsf{COPILOT}}}
\newcommand{\Toracle}{\ensuremath{\mathsf{ORACLE}}}
\newcommand{\Truntime}{\ensuremath{\mathsf{RUNTIME}}}
\newcommand{\Tsolver}{\ensuremath{\mathsf{SOLVER}}}
\newcommand{\Tproof}{\ensuremath{\mathsf{PROOF}}}

% Site and geometry
\newcommand{\Site}{\mathcal{S}}             % semantic site
\newcommand{\Cov}{\mathrm{Cov}}            % covering family
\newcommand{\Ob}{\mathrm{Ob}}              % objects
\newcommand{\Mor}{\mathrm{Mor}}            % morphisms
\newcommand{\res}{\mathrm{res}}            % restriction
\newcommand{\Cech}{\check{C}}              % Čech complex
\newcommand{\Hcoh}[1]{H^{#1}}             % cohomology group

% Descent
\newcommand{\descent}{\mathrm{desc}}
\newcommand{\glue}{\mathrm{glue}}
\newcommand{\overlap}[2]{#1 \cap #2}

% Semantic moves
\newcommand{\VERIFY}{\textsc{Verify}}
\newcommand{\CONSTRUCT}{\textsc{Construct}}
\newcommand{\NORMALIZE}{\textsc{Normalize}}
\newcommand{\GLUE}{\textsc{Glue}}
\newcommand{\RESTRICT}{\textsc{Restrict}}
\newcommand{\TRANSPORT}{\textsc{Transport}}
\newcommand{\REPAIR}{\textsc{Repair}}
\newcommand{\PROMOTE}{\textsc{Promote}}
\newcommand{\CHALLENGE}{\textsc{Challenge}}
\newcommand{\ESCALATE}{\textsc{Escalate}}

% Fragments
\newcommand{\QFLIA}{\texttt{QF\_LIA}}
\newcommand{\QFLRA}{\texttt{QF\_LRA}}
\newcommand{\QFBV}{\texttt{QF\_BV}}
\newcommand{\QFUF}{\texttt{QF\_UF}}

% Misc
\newcommand{\Python}{\textsc{Python}}
\newcommand{\JuGeo}{\textsc{JuGeo}}
\newcommand{\LEAN}{\textsc{Lean}\,4}
\newcommand{\Fstar}{F$^{\star}$}
\newcommand{\Coq}{\textsc{Coq}}
\newcommand{\Dafny}{\textsc{Dafny}}

% Common shortcuts
\newcommand{\ie}{i.e.\@\xspace}
\newcommand{\eg}{e.g.\@\xspace}
\newcommand{\cf}{cf.\@\xspace}
\newcommand{\wrt}{w.r.t.\@\xspace}

% Evaluation shortcuts
\newcommand{\numcases}{300}
\newcommand{\numspec}{100}
\newcommand{\numeq}{100}
\newcommand{\numbug}{100}
\newcommand{\medtime}{6.5\,ms}
\newcommand{\totaltime}{1.949\,s}


% ── Packages ────────────────────────────────────────────────────────────
\usepackage{amsmath,amssymb,amsthm,mathtools}
\usepackage{tikz-cd}
\usepackage{listings}
\usepackage{xcolor}
\usepackage{xspace}
\usepackage{booktabs}
\usepackage{multirow}
\usepackage{enumitem}
\usepackage[expansion=false]{microtype}
\usepackage{hyperref}
\usepackage{cleveref}
% \usepackage{thmtools}  % disabled: not available in this TeX installation
\usepackage{float}
\usepackage{caption}
\usepackage{subcaption}
\usepackage{tabularx}
\usepackage{stmaryrd}       % \llbracket, \rrbracket
\usepackage{bbm}            % \mathbbm{1}

% ── Page geometry ───────────────────────────────────────────────────────
\usepackage[margin=1in]{geometry}

% ── Theorem environments ────────────────────────────────────────────────
\theoremstyle{plain}
\newtheorem{theorem}{Theorem}[section]
\newtheorem{lemma}[theorem]{Lemma}
\newtheorem{proposition}[theorem]{Proposition}
\newtheorem{corollary}[theorem]{Corollary}

\theoremstyle{definition}
\newtheorem{definition}[theorem]{Definition}
\newtheorem{example}[theorem]{Example}
\newtheorem{construction}[theorem]{Construction}

\theoremstyle{remark}
\newtheorem{remark}[theorem]{Remark}
\newtheorem{notation}[theorem]{Notation}

% ── Python listings style ──────────────────────────────────────────────
\definecolor{jugeoBlue}{HTML}{2563EB}
\definecolor{jugeoGreen}{HTML}{059669}
\definecolor{jugeoRed}{HTML}{DC2626}
\definecolor{jugeoPurple}{HTML}{7C3AED}
\definecolor{jugeoGray}{HTML}{6B7280}
\definecolor{jugeoBg}{HTML}{F8FAFC}

\lstdefinestyle{jugeo-python}{
  language=Python,
  basicstyle=\ttfamily\small,
  keywordstyle=\color{jugeoBlue}\bfseries,
  stringstyle=\color{jugeoGreen},
  commentstyle=\color{jugeoGray}\itshape,
  numberstyle=\tiny\color{jugeoGray},
  backgroundcolor=\color{jugeoBg},
  numbers=left,
  numbersep=6pt,
  frame=single,
  framerule=0.4pt,
  rulecolor=\color{jugeoGray!40},
  breaklines=true,
  breakatwhitespace=true,
  showstringspaces=false,
  tabsize=4,
  xleftmargin=1.5em,
  framexleftmargin=1.5em,
  morekeywords={dataclass,frozen,field,Enum,IntEnum,auto,Any,
                Mapping,Sequence,Optional,match,case},
  literate={→}{{$\to$}}1 {←}{{$\leftarrow$}}1
           {≤}{{$\leq$}}1 {≥}{{$\geq$}}1 {≠}{{$\neq$}}1
           {∈}{{$\in$}}1 {∀}{{$\forall$}}1 {∃}{{$\exists$}}1
           {⊕}{{$\oplus$}}1 {⊖}{{$\ominus$}}1,
  escapeinside={(*@}{@*)},
}
\lstset{style=jugeo-python}

\lstdefinestyle{jugeo-output}{
  basicstyle=\ttfamily\small,
  backgroundcolor=\color{jugeoBg},
  frame=single,
  framerule=0.4pt,
  rulecolor=\color{jugeoGray!40},
  breaklines=true,
  showstringspaces=false,
  xleftmargin=1.5em,
  framexleftmargin=0.5em,
  numbers=none,
}

% ── JuGeo-specific macros ──────────────────────────────────────────────

% Judgment 8-tuple
\newcommand{\J}{\mathcal{J}}
\newcommand{\Jud}[8]{(#1,\,#2,\,#3,\,#4,\,#5,\,#6,\,#7,\,#8)}
\newcommand{\JudShort}[1]{\J_{#1}}

% Components
\newcommand{\coord}{\mathsf{c}}            % coordinate
\newcommand{\prop}{\varphi}                 % proposition
\newcommand{\carrier}{\mathcal{A}}          % carrier type
\newcommand{\evid}{\mathcal{E}}             % evidence bundle
\newcommand{\oblig}{\mathcal{O}}            % obligations
\newcommand{\obstr}{\mathcal{B}}            % obstructions
\newcommand{\trust}{\mathcal{T}}            % trust annotation
\newcommand{\prov}{\Pi}                     % provenance

% Trust algebra
\newcommand{\Talg}{\mathfrak{T}}            % trust ordered algebra
\newcommand{\Eadm}{\mathcal{E}_{\mathrm{adm}}}
\newcommand{\tleq}{\preceq}                 % trust ordering
\newcommand{\tjoin}{\oplus}                 % conservative join
\newcommand{\tatten}{\ominus}               % attenuation
\newcommand{\tpromote}{\uparrow_\pi}        % promotion
\newcommand{\tdemote}{\downarrow_\chi}       % demotion

% Trust tiers
\newcommand{\Tcontra}{\ensuremath{\mathsf{CONTRADICTED}}}
\newcommand{\Tunver}{\ensuremath{\mathsf{UNVERIFIED}}}
\newcommand{\Tcopilot}{\ensuremath{\mathsf{COPILOT}}}
\newcommand{\Toracle}{\ensuremath{\mathsf{ORACLE}}}
\newcommand{\Truntime}{\ensuremath{\mathsf{RUNTIME}}}
\newcommand{\Tsolver}{\ensuremath{\mathsf{SOLVER}}}
\newcommand{\Tproof}{\ensuremath{\mathsf{PROOF}}}

% Site and geometry
\newcommand{\Site}{\mathcal{S}}             % semantic site
\newcommand{\Cov}{\mathrm{Cov}}            % covering family
\newcommand{\Ob}{\mathrm{Ob}}              % objects
\newcommand{\Mor}{\mathrm{Mor}}            % morphisms
\newcommand{\res}{\mathrm{res}}            % restriction
\newcommand{\Cech}{\check{C}}              % Čech complex
\newcommand{\Hcoh}[1]{H^{#1}}             % cohomology group

% Descent
\newcommand{\descent}{\mathrm{desc}}
\newcommand{\glue}{\mathrm{glue}}
\newcommand{\overlap}[2]{#1 \cap #2}

% Semantic moves
\newcommand{\VERIFY}{\textsc{Verify}}
\newcommand{\CONSTRUCT}{\textsc{Construct}}
\newcommand{\NORMALIZE}{\textsc{Normalize}}
\newcommand{\GLUE}{\textsc{Glue}}
\newcommand{\RESTRICT}{\textsc{Restrict}}
\newcommand{\TRANSPORT}{\textsc{Transport}}
\newcommand{\REPAIR}{\textsc{Repair}}
\newcommand{\PROMOTE}{\textsc{Promote}}
\newcommand{\CHALLENGE}{\textsc{Challenge}}
\newcommand{\ESCALATE}{\textsc{Escalate}}

% Fragments
\newcommand{\QFLIA}{\texttt{QF\_LIA}}
\newcommand{\QFLRA}{\texttt{QF\_LRA}}
\newcommand{\QFBV}{\texttt{QF\_BV}}
\newcommand{\QFUF}{\texttt{QF\_UF}}

% Misc
\newcommand{\Python}{\textsc{Python}}
\newcommand{\JuGeo}{\textsc{JuGeo}}
\newcommand{\LEAN}{\textsc{Lean}\,4}
\newcommand{\Fstar}{F$^{\star}$}
\newcommand{\Coq}{\textsc{Coq}}
\newcommand{\Dafny}{\textsc{Dafny}}

% Common shortcuts
\newcommand{\ie}{i.e.\@\xspace}
\newcommand{\eg}{e.g.\@\xspace}
\newcommand{\cf}{cf.\@\xspace}
\newcommand{\wrt}{w.r.t.\@\xspace}

% Evaluation shortcuts
\newcommand{\numcases}{300}
\newcommand{\numspec}{100}
\newcommand{\numeq}{100}
\newcommand{\numbug}{100}
\newcommand{\medtime}{6.5\,ms}
\newcommand{\totaltime}{1.949\,s}

% experiment-data.tex — AUTO-GENERATED from real jugeo CLI runs
% DO NOT EDIT — regenerate with: python3 experiments/run_universal_metrics.py
% Generated from 30 programs across 6 domains

\newcommand{\expTotalPrograms}{30}
\newcommand{\expVerified}{30}
\newcommand{\expAccuracy}{100.0\%}
\newcommand{\expCoordsMin}{2}
\newcommand{\expCoordsMax}{10}
\newcommand{\expCoordsMean}{5.2}
\newcommand{\expCoordsMedian}{5.5}
\newcommand{\expPropsMin}{4}
\newcommand{\expPropsMax}{20}
\newcommand{\expPropsMean}{10.4}
\newcommand{\expPropsSum}{311}
\newcommand{\expPropsOkSum}{310}
\newcommand{\expTimeMin}{3.506\,s}
\newcommand{\expTimeMax}{6.714\,s}
\newcommand{\expTimeMean}{4.64\,s}
\newcommand{\expTimeMedian}{4.508\,s}
\newcommand{\expTimeTotal}{139.204\,s}
\newcommand{\expObstructionTotal}{0}
\newcommand{\expHOneZero}{30}
\newcommand{\expDomainCount}{6}
\newcommand{\expTrustCOPILOTSUGGESTED}{30}
\newcommand{\expDomainDsCount}{5}
\newcommand{\expDomainDsVerified}{5}
\newcommand{\expDomainDsAvgCoords}{7.6}
\newcommand{\expDomainDsAvgProps}{15.2}
\newcommand{\expDomainMathCount}{5}
\newcommand{\expDomainMathVerified}{5}
\newcommand{\expDomainMathAvgCoords}{4.8}
\newcommand{\expDomainMathAvgProps}{9.4}
\newcommand{\expDomainSmCount}{5}
\newcommand{\expDomainSmVerified}{5}
\newcommand{\expDomainSmAvgCoords}{7.6}
\newcommand{\expDomainSmAvgProps}{15.2}
\newcommand{\expDomainSortCount}{5}
\newcommand{\expDomainSortVerified}{5}
\newcommand{\expDomainSortAvgCoords}{2.4}
\newcommand{\expDomainSortAvgProps}{4.8}
\newcommand{\expDomainStrCount}{5}
\newcommand{\expDomainStrVerified}{5}
\newcommand{\expDomainStrAvgCoords}{3.2}
\newcommand{\expDomainStrAvgProps}{6.4}
\newcommand{\expDomainWebCount}{5}
\newcommand{\expDomainWebVerified}{5}
\newcommand{\expDomainWebAvgCoords}{5.6}
\newcommand{\expDomainWebAvgProps}{11.2}

% subsystem-data.tex — AUTO-GENERATED from real jugeo subsystem experiments
% DO NOT EDIT — regenerate with: python3 experiments/run_subsystem_experiments.py

\newcommand{\subCliTotal}{10}
\newcommand{\subCliVerified}{9}
\newcommand{\subCliAccuracy}{90.0\%}
\newcommand{\subCliMeanTime}{4.132\,s}
\newcommand{\subCliMinTime}{3.472\,s}
\newcommand{\subCliMaxTime}{5.372\,s}
\newcommand{\subCliTotalCoords}{54}
\newcommand{\subCliTotalProps}{107}
\newcommand{\subCliTotalPropsOk}{106}
\newcommand{\subSiteCalls}{90}
\newcommand{\subSiteSuccessful}{69}
\newcommand{\subSiteSuccessRate}{76.7\%}
\newcommand{\subSiteMeanTime}{0.0001\,s}
\newcommand{\subSiteTotalTime}{0.005\,s}
\newcommand{\subSpecificationsatisfactionSmallTime}{0.0003\,s}
\newcommand{\subSpecificationsatisfactionMediumTime}{0.0001\,s}
\newcommand{\subSpecificationsatisfactionLargeTime}{0.0001\,s}
\newcommand{\subInhabitantfleetSmallTime}{0.0\,s}
\newcommand{\subInhabitantfleetMediumTime}{0.0\,s}
\newcommand{\subInhabitantfleetLargeTime}{0.0\,s}
\newcommand{\subTheoremecologySmallTime}{0.0\,s}
\newcommand{\subTheoremecologyMediumTime}{0.0\,s}
\newcommand{\subTheoremecologyLargeTime}{0.0\,s}
\newcommand{\subStatespaceexplorationSmallTime}{0.0\,s}
\newcommand{\subStatespaceexplorationMediumTime}{0.0\,s}
\newcommand{\subStatespaceexplorationLargeTime}{0.0\,s}
\newcommand{\subRepairsemanticsSmallTime}{0.0\,s}
\newcommand{\subRepairsemanticsMediumTime}{0.0\,s}
\newcommand{\subRepairsemanticsLargeTime}{0.0\,s}
\newcommand{\subReplaygluingSmallTime}{0.0\,s}
\newcommand{\subReplaygluingMediumTime}{0.0\,s}
\newcommand{\subReplaygluingLargeTime}{0.0\,s}
\newcommand{\subSemanticfuturesSmallTime}{0.0\,s}
\newcommand{\subSemanticfuturesMediumTime}{0.0\,s}
\newcommand{\subSemanticfuturesLargeTime}{0.0\,s}
\newcommand{\subHypercovertreatySmallTime}{0.0\,s}
\newcommand{\subHypercovertreatyMediumTime}{0.0\,s}
\newcommand{\subHypercovertreatyLargeTime}{0.0\,s}
\newcommand{\subEncodeforsolverSmallTime}{0.0026\,s}
\newcommand{\subEncodeforsolverMediumTime}{0.0002\,s}
\newcommand{\subEncodeforsolverLargeTime}{0.0001\,s}
\newcommand{\subInterfaceroutingSmallTime}{0.0\,s}
\newcommand{\subInterfaceroutingMediumTime}{0.0\,s}
\newcommand{\subInterfaceroutingLargeTime}{0.0\,s}
\newcommand{\subOrchestrateverificationSmallTime}{0.0002\,s}
\newcommand{\subOrchestrateverificationMediumTime}{0.0\,s}
\newcommand{\subOrchestrateverificationLargeTime}{0.0\,s}
\newcommand{\subRunfulldescentSmallTime}{0.0\,s}
\newcommand{\subRunfulldescentMediumTime}{0.0\,s}
\newcommand{\subRunfulldescentLargeTime}{0.0\,s}
\newcommand{\subKernellifecycleSmallTime}{0.0\,s}
\newcommand{\subKernellifecycleMediumTime}{0.0\,s}
\newcommand{\subKernellifecycleLargeTime}{0.0\,s}
\newcommand{\subDiscoverypipelineSmallTime}{0.0\,s}
\newcommand{\subDiscoverypipelineMediumTime}{0.0\,s}
\newcommand{\subDiscoverypipelineLargeTime}{0.0\,s}
\newcommand{\subAnalogytransportSmallTime}{0.0\,s}
\newcommand{\subAnalogytransportMediumTime}{0.0\,s}
\newcommand{\subAnalogytransportLargeTime}{0.0\,s}
\newcommand{\subFormalcoresiteSmallTime}{0.0001\,s}
\newcommand{\subFormalcoresiteMediumTime}{0.0\,s}
\newcommand{\subFormalcoresiteLargeTime}{0.0\,s}
\newcommand{\subProblematlasSmallTime}{0.0\,s}
\newcommand{\subProblematlasMediumTime}{0.0\,s}
\newcommand{\subProblematlasLargeTime}{0.0\,s}
\newcommand{\subPublicalignmentSmallTime}{0.0\,s}
\newcommand{\subPublicalignmentMediumTime}{0.0\,s}
\newcommand{\subPublicalignmentLargeTime}{0.0\,s}
\newcommand{\subRegimebootstrappingSmallTime}{0.0\,s}
\newcommand{\subRegimebootstrappingMediumTime}{0.0\,s}
\newcommand{\subRegimebootstrappingLargeTime}{0.0\,s}
\newcommand{\subRelationalrefinementSmallTime}{0.0\,s}
\newcommand{\subRelationalrefinementMediumTime}{0.0\,s}
\newcommand{\subRelationalrefinementLargeTime}{0.0\,s}
\newcommand{\subSemanticclosureSmallTime}{0.0\,s}
\newcommand{\subSemanticclosureMediumTime}{0.0\,s}
\newcommand{\subSemanticclosureLargeTime}{0.0\,s}
\newcommand{\subTheoremeconomicsSmallTime}{0.0001\,s}
\newcommand{\subTheoremeconomicsMediumTime}{0.0\,s}
\newcommand{\subTheoremeconomicsLargeTime}{0.0\,s}
\newcommand{\subBenchmarksuiteSmallTime}{0.0\,s}
\newcommand{\subBenchmarksuiteMediumTime}{0.0\,s}
\newcommand{\subBenchmarksuiteLargeTime}{0.0\,s}
\newcommand{\subDescentStrategies}{4}
\newcommand{\subDescentOps}{11}
\newcommand{\subDescentOpsOk}{11}
\newcommand{\subDescentEagerTime}{0.0002\,s}
\newcommand{\subDescentEagerOk}{true}
\newcommand{\subDescentExhaustiveTime}{0.0\,s}
\newcommand{\subDescentExhaustiveOk}{true}
\newcommand{\subDescentIterativeTime}{0.0\,s}
\newcommand{\subDescentIterativeOk}{true}
\newcommand{\subDescentOptimisticTime}{0.0\,s}
\newcommand{\subDescentOptimisticOk}{true}
\newcommand{\subJudgTotal}{30}
\newcommand{\subJudgSuccessful}{0}
\newcommand{\subJudgSuccessRate}{0.0\%}
\newcommand{\subJudgMeanTimeUs}{7.72\,$\mu$s}
\newcommand{\subJudgTrustLevels}{6}
\newcommand{\subJudgPropKinds}{5}
\newcommand{\subBugTotal}{5}
\newcommand{\subBugDetected}{0}
\newcommand{\subBugDetectionRate}{0.0\%}
\newcommand{\subBugFalsePositiveRate}{0.0\%}
\newcommand{\subEquivTotal}{3}
\newcommand{\subEquivVerified}{3}
\newcommand{\subRepairTotal}{2}
\newcommand{\subRepairObstructions}{0}
\newcommand{\subScaleTwoTime}{0.0001\,s}
\newcommand{\subScaleFourTime}{0.0\,s}
\newcommand{\subScaleEightTime}{0.0\,s}
\newcommand{\subScaleTwelveTime}{0.0\,s}
\newcommand{\subScaleSixteenTime}{0.0\,s}
\newcommand{\subScaleTwentyTime}{0.0\,s}
\newcommand{\subTotalExperimentTime}{95.6\,s}

% comprehensive-data.tex — AUTO-GENERATED from real jugeo experiments
% DO NOT EDIT — regenerate with: python3 experiments/run_comprehensive_experiments.py
% Generated: 2026-03-23 09:19:06

% --- Size scaling metrics ---
\newcommand{\compTotalPrograms}{18}
\newcommand{\compTotalVerified}{18}
\newcommand{\compOverallAccuracy}{100.0\%}
\newcommand{\compTinyCount}{5}
\newcommand{\compTinyVerified}{5}
\newcommand{\compTinyMeanCoords}{3}
\newcommand{\compTinyMeanProps}{6}
\newcommand{\compTinyMeanTime}{4.95\,s}
\newcommand{\compTinyMeanLines}{5}
\newcommand{\compTinyMinTime}{4.45\,s}
\newcommand{\compTinyMaxTime}{5.51\,s}
\newcommand{\compSmallCount}{5}
\newcommand{\compSmallVerified}{5}
\newcommand{\compSmallMeanCoords}{3.6}
\newcommand{\compSmallMeanProps}{7.2}
\newcommand{\compSmallMeanTime}{4.85\,s}
\newcommand{\compSmallMeanLines}{16.0}
\newcommand{\compSmallMinTime}{4.30\,s}
\newcommand{\compSmallMaxTime}{5.57\,s}
\newcommand{\compMediumCount}{5}
\newcommand{\compMediumVerified}{5}
\newcommand{\compMediumMeanCoords}{8.6}
\newcommand{\compMediumMeanProps}{17.2}
\newcommand{\compMediumMeanTime}{4.47\,s}
\newcommand{\compMediumMeanLines}{45.0}
\newcommand{\compMediumMinTime}{4.26\,s}
\newcommand{\compMediumMaxTime}{4.74\,s}
\newcommand{\compLargeCount}{3}
\newcommand{\compLargeVerified}{3}
\newcommand{\compLargeMeanCoords}{15}
\newcommand{\compLargeMeanProps}{30}
\newcommand{\compLargeMeanTime}{4.31\,s}
\newcommand{\compLargeMeanLines}{79.0}
\newcommand{\compLargeMinTime}{4.27\,s}
\newcommand{\compLargeMaxTime}{4.38\,s}
% --- Aggregate timing ---
\newcommand{\compTimeMean}{4.68\,s}
\newcommand{\compTimeMedian}{4.50\,s}
\newcommand{\compTimeMin}{4.265\,s}
\newcommand{\compTimeMax}{5.571\,s}
\newcommand{\compTimeTotal}{84.3\,s}
\newcommand{\compTimeStdev}{0.45\,s}
% --- Aggregate structure ---
\newcommand{\compCoordsMean}{6.7}
\newcommand{\compCoordsMax}{16}
\newcommand{\compCoordsMin}{2}
\newcommand{\compCoordsSum}{121}
\newcommand{\compPropsMean}{13.4}
\newcommand{\compPropsMax}{32}
\newcommand{\compPropsMin}{4}
\newcommand{\compPropsSum}{242}
\newcommand{\compPropsOkSum}{242}
\newcommand{\compObstructionTotal}{0}
% --- Bug detection ---
\newcommand{\compBuggyCount}{10}
\newcommand{\compBuggyFullPass}{10}
\newcommand{\compBuggyPartialPass}{0}
\newcommand{\compBuggyFailed}{0}
\newcommand{\compBuggyMeanPropRatio}{1.00}
\newcommand{\compCorrectMeanPropRatio}{1.00}
\newcommand{\compBuggyMeanTime}{5.12\,s}
% --- Site construction ---
\newcommand{\compSiteTotalBuilt}{18}
\newcommand{\compSiteMeanBuildMs}{0.05}
\newcommand{\compSiteMaxBuildMs}{0.14}
\newcommand{\compSiteMeanCoords}{6.5}
\newcommand{\compSiteMeanMorphisms}{14.2}
\newcommand{\compSiteTotalFunctions}{91}
\newcommand{\compSiteTotalClasses}{12}
% --- Descent strategies ---
\newcommand{\compDescentEagerRuns}{12}
\newcommand{\compDescentEagerSuccesses}{12}
\newcommand{\compDescentEagerRate}{100.0\%}
\newcommand{\compDescentEagerMeanMs}{0.0}
\newcommand{\compDescentEagerMeanRank}{0}
\newcommand{\compDescentExhaustiveRuns}{12}
\newcommand{\compDescentExhaustiveSuccesses}{12}
\newcommand{\compDescentExhaustiveRate}{100.0\%}
\newcommand{\compDescentExhaustiveMeanMs}{0.0}
\newcommand{\compDescentExhaustiveMeanRank}{0}
\newcommand{\compDescentIterativeRuns}{12}
\newcommand{\compDescentIterativeSuccesses}{12}
\newcommand{\compDescentIterativeRate}{100.0\%}
\newcommand{\compDescentIterativeMeanMs}{0.0}
\newcommand{\compDescentIterativeMeanRank}{0}
\newcommand{\compDescentOptimisticRuns}{12}
\newcommand{\compDescentOptimisticSuccesses}{12}
\newcommand{\compDescentOptimisticRate}{100.0\%}
\newcommand{\compDescentOptimisticMeanMs}{0.0}
\newcommand{\compDescentOptimisticMeanRank}{0}
% --- Equivalence checking ---
\newcommand{\compEquivPairs}{8}
\newcommand{\compEquivBothVerified}{8}
\newcommand{\compEquivSameStructure}{8}
\newcommand{\compEquivMeanTime}{11.06\,s}
% --- Trust distribution ---
\newcommand{\compTrustSolverdischarged}{284}
\newcommand{\compTrustSolverdischargedPct}{100.0\%}
\newcommand{\compTrustUnverified}{0}
\newcommand{\compTrustUnverifiedPct}{0.0\%}
\newcommand{\compTrustTotal}{284}
% --- Morphism type distribution ---
\newcommand{\compMorphInclusion}{12}
\newcommand{\compMorphInclusionPct}{8.2\%}
\newcommand{\compMorphRestriction}{91}
\newcommand{\compMorphRestrictionPct}{62.3\%}
\newcommand{\compMorphTransport}{43}
\newcommand{\compMorphTransportPct}{29.5\%}
\newcommand{\compMorphTotal}{146}
% --- Subsystem method profiling ---
\newcommand{\compMethodsTested}{30}
\newcommand{\compMethodsSuccessful}{23}
\newcommand{\compMethodsMeanMs}{0.02}
\newcommand{\compProfAnalogyTransportMeanMs}{0.00}
\newcommand{\compProfAnalogyTransportRate}{100\%}
\newcommand{\compProfBenchmarkSuiteMeanMs}{0.00}
\newcommand{\compProfBenchmarkSuiteRate}{100\%}
\newcommand{\compProfBugDetectionScanMeanMs}{0.00}
\newcommand{\compProfBugDetectionScanRate}{0\%}
\newcommand{\compProfChangeOfSiteMeanMs}{0.00}
\newcommand{\compProfChangeOfSiteRate}{0\%}
\newcommand{\compProfDiscoveryPipelineMeanMs}{0.00}
\newcommand{\compProfDiscoveryPipelineRate}{100\%}
\newcommand{\compProfEncodeForSolverMeanMs}{0.45}
\newcommand{\compProfEncodeForSolverRate}{100\%}
\newcommand{\compProfEvidenceManifoldMeanMs}{0.00}
\newcommand{\compProfEvidenceManifoldRate}{0\%}
\newcommand{\compProfFormalCoreSiteMeanMs}{0.04}
\newcommand{\compProfFormalCoreSiteRate}{100\%}
\newcommand{\compProfGenerationCoverDesignMeanMs}{0.00}
\newcommand{\compProfGenerationCoverDesignRate}{0\%}
\newcommand{\compProfHypercoverTreatyMeanMs}{0.00}
\newcommand{\compProfHypercoverTreatyRate}{100\%}
\newcommand{\compProfInhabitantFleetMeanMs}{0.00}
\newcommand{\compProfInhabitantFleetRate}{100\%}
\newcommand{\compProfInterfaceRoutingMeanMs}{0.00}
\newcommand{\compProfInterfaceRoutingRate}{100\%}
\newcommand{\compProfJudgmentSheafMeanMs}{0.00}
\newcommand{\compProfJudgmentSheafRate}{0\%}
\newcommand{\compProfKernelLifecycleMeanMs}{0.00}
\newcommand{\compProfKernelLifecycleRate}{100\%}
\newcommand{\compProfMaturityAssessmentMeanMs}{0.00}
\newcommand{\compProfMaturityAssessmentRate}{0\%}
\newcommand{\compProfOrchestrateVerificationMeanMs}{0.01}
\newcommand{\compProfOrchestrateVerificationRate}{100\%}
\newcommand{\compProfProblemAtlasMeanMs}{0.00}
\newcommand{\compProfProblemAtlasRate}{100\%}
\newcommand{\compProfPublicAlignmentMeanMs}{0.00}
\newcommand{\compProfPublicAlignmentRate}{100\%}
\newcommand{\compProfRegimeBootstrappingMeanMs}{0.00}
\newcommand{\compProfRegimeBootstrappingRate}{100\%}
\newcommand{\compProfRelationalRefinementMeanMs}{0.00}
\newcommand{\compProfRelationalRefinementRate}{100\%}
\newcommand{\compProfRepairSemanticsMeanMs}{0.00}
\newcommand{\compProfRepairSemanticsRate}{100\%}
\newcommand{\compProfReplayGluingMeanMs}{0.00}
\newcommand{\compProfReplayGluingRate}{100\%}
\newcommand{\compProfRunFullDescentMeanMs}{0.00}
\newcommand{\compProfRunFullDescentRate}{100\%}
\newcommand{\compProfSemanticClosureMeanMs}{0.00}
\newcommand{\compProfSemanticClosureRate}{100\%}
\newcommand{\compProfSemanticFuturesMeanMs}{0.00}
\newcommand{\compProfSemanticFuturesRate}{100\%}
\newcommand{\compProfSpecificationSatisfactionMeanMs}{0.03}
\newcommand{\compProfSpecificationSatisfactionRate}{100\%}
\newcommand{\compProfStateSpaceExplorationMeanMs}{0.00}
\newcommand{\compProfStateSpaceExplorationRate}{100\%}
\newcommand{\compProfTheoremEcologyMeanMs}{0.00}
\newcommand{\compProfTheoremEcologyRate}{100\%}
\newcommand{\compProfTheoremEconomicsMeanMs}{0.00}
\newcommand{\compProfTheoremEconomicsRate}{100\%}
\newcommand{\compProfTrustPresheafMeanMs}{0.00}
\newcommand{\compProfTrustPresheafRate}{0\%}

% supplementary-data.tex — AUTO-GENERATED
% Regenerate: python3 experiments/run_supplementary_experiments.py
% Generated: 2026-03-23 09:57:40

\newcommand{\suppDomainSortAvgTime}{3.59\,s}
\newcommand{\suppDomainSortMinTime}{3.18\,s}
\newcommand{\suppDomainSortMaxTime}{4.00\,s}
\newcommand{\suppDomainDsAvgTime}{3.41\,s}
\newcommand{\suppDomainDsMinTime}{3.27\,s}
\newcommand{\suppDomainDsMaxTime}{3.75\,s}
\newcommand{\suppDomainMathAvgTime}{3.76\,s}
\newcommand{\suppDomainMathMinTime}{3.15\,s}
\newcommand{\suppDomainMathMaxTime}{4.05\,s}
\newcommand{\suppDomainStrAvgTime}{3.27\,s}
\newcommand{\suppDomainStrMinTime}{3.05\,s}
\newcommand{\suppDomainStrMaxTime}{3.47\,s}
\newcommand{\suppDomainWebAvgTime}{3.33\,s}
\newcommand{\suppDomainWebMinTime}{3.25\,s}
\newcommand{\suppDomainWebMaxTime}{3.47\,s}
\newcommand{\suppDomainSmAvgTime}{3.81\,s}
\newcommand{\suppDomainSmMinTime}{3.41\,s}
\newcommand{\suppDomainSmMaxTime}{4.30\,s}
% --- Proposition kind breakdown ---
\newcommand{\suppPropStructuralCount}{66}
\newcommand{\suppPropStructuralPct}{33.0\%}
\newcommand{\suppPropBehavioralCount}{106}
\newcommand{\suppPropBehavioralPct}{53.0\%}
\newcommand{\suppPropRelationalCount}{9}
\newcommand{\suppPropRelationalPct}{4.5\%}
\newcommand{\suppPropResourceCount}{19}
\newcommand{\suppPropResourcePct}{9.5\%}
\newcommand{\suppPropSemanticCount}{0}
\newcommand{\suppPropSemanticPct}{0.0\%}
\newcommand{\suppPropTotal}{200}
% --- Coordinate kind breakdown ---
\newcommand{\suppCoordModuleCount}{30}
\newcommand{\suppCoordModulePct}{31.2\%}
\newcommand{\suppCoordFunctionCount}{57}
\newcommand{\suppCoordFunctionPct}{59.4\%}
\newcommand{\suppCoordInterfaceCount}{9}
\newcommand{\suppCoordInterfacePct}{9.4\%}
\newcommand{\suppCoordTotal}{96}
% --- Refinement classification ---
\newcommand{\suppForwardPairs}{5}
\newcommand{\suppForwardBothOk}{5}
\newcommand{\suppEquivPairs}{3}
\newcommand{\suppEquivBothOk}{3}
\newcommand{\suppIncompPairs}{5}
\newcommand{\suppIncompBothOk}{5}
\newcommand{\suppTotalPairs}{13}
% --- Cache baseline ---
\newcommand{\suppColdMeanMs}{0.663}
\newcommand{\suppWarmMeanMs}{0.019}
\newcommand{\suppCacheSpeedup}{35.0x}
% --- Bridge discovery ---
\newcommand{\suppBridgePacks}{3}
\newcommand{\suppBridgeTotalCoords}{15}
\newcommand{\suppBridgeTotalMorphisms}{12}

% data-paper64.tex — AUTO-GENERATED by exp64_team_workflow.py
% DO NOT EDIT — regenerate with: python3 experiments/exp64_team_workflow.py
% Generated from 8 team programs

\newcommand{\ppLXIVteamCount}{8}
\newcommand{\ppLXIVteamsOk}{8}
\newcommand{\ppLXIVfullVerified}{8}
\newcommand{\ppLXIVfullVerifiedPct}{100.0\%}
\newcommand{\ppLXIVmergeConsistent}{8}
\newcommand{\ppLXIVmergeConsistentPct}{100.0\%}

\newcommand{\ppLXIVcoordsMean}{8.5}
\newcommand{\ppLXIVpropsMean}{17}
\newcommand{\ppLXIVoverlapMean}{2}

\newcommand{\ppLXIVtimeMean}{21.17\,s}
\newcommand{\ppLXIVtimeTotal}{169.35\,s}
\newcommand{\ppLXIVtimeMin}{17.718\,s}
\newcommand{\ppLXIVtimeMax}{28.669\,s}

% Per-team results
\newcommand{\ppLXIVteamauthmoduleCoords}{11}
\newcommand{\ppLXIVteamauthmoduleFullProps}{22}
\newcommand{\ppLXIVteamauthmoduleDevAProps}{14}
\newcommand{\ppLXIVteamauthmoduleDevBProps}{10}
\newcommand{\ppLXIVteamauthmoduleOverlap}{2}
\newcommand{\ppLXIVteamauthmoduleMerge}{Yes}
\newcommand{\ppLXIVteamauthmoduleTime}{28.669\,s}
\newcommand{\ppLXIVteamdatapipelineCoords}{6}
\newcommand{\ppLXIVteamdatapipelineFullProps}{12}
\newcommand{\ppLXIVteamdatapipelineDevAProps}{6}
\newcommand{\ppLXIVteamdatapipelineDevBProps}{8}
\newcommand{\ppLXIVteamdatapipelineOverlap}{2}
\newcommand{\ppLXIVteamdatapipelineMerge}{Yes}
\newcommand{\ppLXIVteamdatapipelineTime}{19.978\,s}
\newcommand{\ppLXIVteamwebhandlerCoords}{9}
\newcommand{\ppLXIVteamwebhandlerFullProps}{18}
\newcommand{\ppLXIVteamwebhandlerDevAProps}{10}
\newcommand{\ppLXIVteamwebhandlerDevBProps}{10}
\newcommand{\ppLXIVteamwebhandlerOverlap}{2}
\newcommand{\ppLXIVteamwebhandlerMerge}{Yes}
\newcommand{\ppLXIVteamwebhandlerTime}{21.326\,s}
\newcommand{\ppLXIVteamcachesystemCoords}{10}
\newcommand{\ppLXIVteamcachesystemFullProps}{20}
\newcommand{\ppLXIVteamcachesystemDevAProps}{16}
\newcommand{\ppLXIVteamcachesystemDevBProps}{6}
\newcommand{\ppLXIVteamcachesystemOverlap}{2}
\newcommand{\ppLXIVteamcachesystemMerge}{Yes}
\newcommand{\ppLXIVteamcachesystemTime}{18.422\,s}
\newcommand{\ppLXIVteameventsystemCoords}{10}
\newcommand{\ppLXIVteameventsystemFullProps}{20}
\newcommand{\ppLXIVteameventsystemDevAProps}{12}
\newcommand{\ppLXIVteameventsystemDevBProps}{10}
\newcommand{\ppLXIVteameventsystemOverlap}{2}
\newcommand{\ppLXIVteameventsystemMerge}{Yes}
\newcommand{\ppLXIVteameventsystemTime}{17.718\,s}
\newcommand{\ppLXIVteammathlibraryCoords}{6}
\newcommand{\ppLXIVteammathlibraryFullProps}{12}
\newcommand{\ppLXIVteammathlibraryDevAProps}{6}
\newcommand{\ppLXIVteammathlibraryDevBProps}{8}
\newcommand{\ppLXIVteammathlibraryOverlap}{2}
\newcommand{\ppLXIVteammathlibraryMerge}{Yes}
\newcommand{\ppLXIVteammathlibraryTime}{18.703\,s}
\newcommand{\ppLXIVteamvalidatorCoords}{4}
\newcommand{\ppLXIVteamvalidatorFullProps}{8}
\newcommand{\ppLXIVteamvalidatorDevAProps}{4}
\newcommand{\ppLXIVteamvalidatorDevBProps}{6}
\newcommand{\ppLXIVteamvalidatorOverlap}{2}
\newcommand{\ppLXIVteamvalidatorMerge}{Yes}
\newcommand{\ppLXIVteamvalidatorTime}{21.364\,s}
\newcommand{\ppLXIVteamconfigsystemCoords}{12}
\newcommand{\ppLXIVteamconfigsystemFullProps}{24}
\newcommand{\ppLXIVteamconfigsystemDevAProps}{16}
\newcommand{\ppLXIVteamconfigsystemDevBProps}{10}
\newcommand{\ppLXIVteamconfigsystemOverlap}{2}
\newcommand{\ppLXIVteamconfigsystemMerge}{Yes}
\newcommand{\ppLXIVteamconfigsystemTime}{23.166\,s}


\usepackage{array}
\usepackage{tikz}
\usetikzlibrary{arrows.meta,positioning,shapes.geometric,fit,calc}
\usepackage{algorithm}
\usepackage{algorithmic}

% Paper-specific macros
\newcommand{\TeamSite}{\textsc{TeamSite}}
\newcommand{\AuthGraph}{\textsc{AuthGraph}}
\newcommand{\MergeEngine}{\textsc{MergeEngine}}
\newcommand{\DelegToken}{\mathsf{DELEGATED}}
\newcommand{\OwnerToken}{\mathsf{OWNER}}
\newcommand{\ReviewToken}{\mathsf{REVIEWER}}
\newcommand{\MergeOp}{\texttt{merge\_judgments}}
\newcommand{\DelegOp}{\texttt{delegate\_authority}}
\newcommand{\ConflictRes}{\texttt{resolve\_conflicts}}
\newcommand{\jurisd}{\mathrm{jurisd}}
\newcommand{\auth}{\mathrm{auth}}
\newcommand{\Pre}{\mathrm{Pre}}
\newcommand{\Sh}{\mathrm{Sh}}
\newcommand{\colim}{\mathrm{colim}}
\newcommand{\id}{\mathrm{id}}
\newcommand{\op}{\mathrm{op}}
\newcommand{\Set}{\mathbf{Set}}
\newcommand{\Ab}{\mathbf{Ab}}
\newcommand{\Cat}{\mathbf{Cat}}
\newcommand{\DelegCat}{\mathsf{Deleg}}
\newcommand{\JurisCat}{\mathsf{Juris}}
\newcommand{\MergeFun}{\mathcal{M}}
\newcommand{\TeamCech}{\check{C}_{\mathrm{team}}}
\newcommand{\nerve}{\mathrm{N}}

\title{Multi-Developer Workflows with Jurisdiction and\\Authority Delegation in \JuGeo}
\author{JuGeo Development Team}
\date{\today}

\begin{document}
\maketitle

\begin{abstract}
Modern software projects require multiple developers to collaborate on shared
codebases while maintaining formal verification guarantees. We present a
framework for multi-developer workflows within the \JuGeo{} judgment-geometric
verification system, where each developer operates within a defined
\emph{jurisdiction}---a sub-site of the semantic site---and may delegate
verification authority to teammates via structured trust morphisms. The
empirical part of this paper is limited to the curated team-workflow corpus:
across \ppLXIVteamCount{} team programs, all merged runs remained consistent,
with mean site size \ppLXIVcoordsMean{} coordinates, mean
\ppLXIVpropsMean{} propositions, mean overlap size \ppLXIVoverlapMean{}, and
mean runtime \ppLXIVtimeMean{}. We therefore present the measurements as
descriptive JuGeo-backed data for the benchmark rather than a claim about all
real-world team workflows.
\end{abstract}

% ------------------------------------------------------------------
\section{Introduction}\label{sec:intro}
% ------------------------------------------------------------------

Formal verification of production software is rarely a solo endeavor.
Development teams partition responsibility across modules, features, and
architectural layers. Yet most verification frameworks treat the codebase as a
monolithic artifact owned by a single prover. This mismatch creates friction:
developers must serialize verification effort, duplicate context, or abandon
formal guarantees at module boundaries.

The problem is fundamentally one of \emph{distributed coherence}: how can
independently produced verification artifacts be combined into a consistent
whole? Classical approaches to modular verification---such as assume-guarantee
reasoning~\cite{henzinger1998you} or compositional model
checking~\cite{clarke2000counterexample}---address the technical question of
composing proofs but neglect the organizational dimension. They do not model
\emph{who} is authorized to verify \emph{what}, how authority flows between
team members, or how concurrent verification efforts are reconciled when
developers modify overlapping scopes.

These organizational concerns are not merely administrative. Without formal
treatment, the social protocols of team verification---who verifies what, how
authority flows, how concurrent results are reconciled---remain ad hoc,
undermining the rigor that verification is meant to provide.

\JuGeo{} offers a natural resolution. The judgment 8-tuple
$\Jud{\coord}{\prop}{\carrier}{\evid}{\oblig}{\obstr}{\trust}{\prov}$ already
encodes provenance information $\prov$ that tracks \emph{who} produced a
judgment and \emph{how}. The semantic site $\Site$ organizes code into objects
(scopes) connected by morphisms (restrictions, inclusions, transports). In this
paper, we show how to partition $\Site$ into \emph{jurisdictions}---sub-sites
assigned to individual developers---and define \emph{authority delegation
morphisms} that allow one developer to act on behalf of another within
controlled bounds.

The key insight is that the sheaf-theoretic machinery already provides the
correct framework. A jurisdiction cover
$\{\Site_d\}_{d \in D}$ is a covering family in the Grothendieck topology, and
the gluing axiom guarantees that if developers independently verify their
jurisdictions and agree on overlaps, the result is a \emph{provably sound}
global verification by the universal property of the sheaf.

We further connect team conflicts to \v{C}ech cohomology:
$\Hcoh{1}(\Site, \J)$ classifies collections of local judgments that fail to
glue. When $\Hcoh{1}(\Site, \J) = 0$, every compatible family extends globally.
Our experiments confirm this vanishing across all \expHOneZero{} programs.

Our contributions are:
\begin{enumerate}[leftmargin=*]
  \item A formal model of jurisdictions as sub-sites with covering semantics,
        including a jurisdiction lattice with refinement ordering, ensuring that
        the union of all jurisdictions covers the full site
        (Section~\ref{sec:jurisdictions}).
  \item A delegation protocol built on trust morphisms with attenuation,
        formalized as a category $\DelegCat$ of delegation morphisms,
        preventing unbounded privilege escalation
        (Section~\ref{sec:delegation}).
  \item A merge algorithm for concurrent judgment streams that detects and
        resolves conflicts using obstruction theory, with a proof that the
        merge functor $\MergeFun$ preserves the sheaf condition
        (Section~\ref{sec:merge}).
  \item A cohomological analysis of team conflicts via the team \v{C}ech
        complex $\TeamCech$, connecting $\Hcoh{1} = 0$ to successful global
        merge and establishing a spectral sequence for hierarchical teams
        (Section~\ref{sec:cohomology}).
  \item An evaluation on \expTotalPrograms{} programs across \expDomainCount{}
        domains showing zero merge-induced obstructions, with detailed
        per-domain timing, proposition breakdowns, and scaling analysis
        (Section~\ref{sec:evaluation}).
\end{enumerate}

% ------------------------------------------------------------------
\section{Background}\label{sec:background}
% ------------------------------------------------------------------

\subsection{Judgment-Geometric Verification}

\JuGeo{} verifies programs by constructing a semantic site $\Site$ whose objects
represent program scopes (functions, classes, modules) and whose morphisms
capture structural relationships. For each object $U \in \Ob(\Site)$, the
system produces judgments $\J$ whose propositions $\prop$ encode correctness
properties. Verification succeeds when a global section of the judgment sheaf
exists---that is, when local judgments glue consistently across overlaps.

The trust algebra $\Talg$ orders evidence quality from $\Tcontra$ (contradicted)
through $\Tunver$, $\Tcopilot$, $\Toracle$, $\Truntime$, $\Tsolver$, to
$\Tproof$ (formally proved). The conservative join $\tjoin$ combines evidence
from independent sources, while attenuation $\tatten$ reduces trust when
crossing jurisdiction boundaries.

\subsection{Categorical Foundations of Team Verification}\label{sec:cat-foundations}

We briefly recall the categorical machinery underlying the sheaf-theoretic
approach. A \emph{Grothendieck topology} $\tau$ on a category $\mathcal{C}$
assigns to each object $U \in \Ob(\mathcal{C})$ a collection
$\tau(U)$ of \emph{covering sieves}---families of morphisms into $U$ that are
``jointly surjective'' in the generalized sense.

\begin{definition}[Grothendieck Topology on $\Site$]
The semantic site $(\Site, \tau)$ is equipped with the topology $\tau$ where a
family $\{f_i : U_i \to U\}_{i \in I}$ is a covering family if and only if:
\begin{enumerate}
  \item \textbf{Scope coverage:} $\bigcup_{i \in I} \mathrm{Im}(f_i) = U$.
  \item \textbf{Pullback stability:} For any $g : V \to U$, the pullback
        family $\{U_i \times_U V \to V\}$ is covering.
  \item \textbf{Transitivity:} Compositions of covers are covers.
\end{enumerate}
\end{definition}

\begin{definition}[Presheaf and Sheaf on $\Site$]
A \emph{presheaf} on $\Site$ is a functor $F : \Site^{\op} \to \Set$. It is a
\emph{sheaf} if for every covering family $\{U_i \to U\}_{i \in I}$, the
diagram
\[
  F(U) \xrightarrow{\;\;\prod \res_{U_i}\;\;}
  \prod_{i \in I} F(U_i)
  \rightrightarrows
  \prod_{i,j \in I} F(U_i \times_U U_j)
\]
is an equalizer: compatible local sections extend uniquely to a global section.
\end{definition}

In the \JuGeo{} context, the judgment presheaf $\J : \Site^{\op} \to \Set$ sends
each scope $U$ to the set of judgments $\J(U)$ and each morphism
$f : V \to U$ to the restriction map $\res_f : \J(U) \to \J(V)$. The sheaf
condition for $\J$ is precisely the assertion that locally verified code
composes into a globally verified program.

\begin{definition}[\v{C}ech Cohomology]
Given a covering family $\mathfrak{U} = \{U_i\}_{i \in I}$ of $U$ and a sheaf
$\mathcal{F}$ of abelian groups on $\Site$, the \emph{\v{C}ech complex}
$\Cech^\bullet(\mathfrak{U}, \mathcal{F})$ is defined by
\[
  \Cech^p(\mathfrak{U}, \mathcal{F}) =
  \prod_{i_0 < \cdots < i_p} \mathcal{F}(U_{i_0} \cap \cdots \cap U_{i_p}),
\]
with coboundary maps $\delta^p : \Cech^p \to \Cech^{p+1}$ given by the
alternating sum of restriction maps. The \v{C}ech cohomology groups are
$\Hcoh{p}(\mathfrak{U}, \mathcal{F}) = \ker \delta^p / \mathrm{im}\,\delta^{p-1}$.
\end{definition}

The zeroth cohomology $\Hcoh{0}(\mathfrak{U}, \J)$ recovers global
sections---successfully glued judgments. The first $\Hcoh{1}(\mathfrak{U}, \J)$
classifies obstructions to gluing. When $\Hcoh{1} = 0$, every compatible
family glues, the condition for team merge soundness.

\subsection{Authority Delegation (Paper 19)}

Paper~19 introduced authority delegation as a trust morphism
$\delta : \trust_A \to \trust_B$ between developer trust contexts. A delegation
token $\DelegToken$ records the delegator, delegatee, scope restriction, and
expiration. We extend this model with team-wide coordination.

The key insight from Paper~19 is that delegation is a \emph{morphism in the
trust algebra} $\Talg$: the attenuation $\tatten$ acts as a functor preserving
$\tleq$ while strictly decreasing trust levels, enabling compositional chaining.

\subsection{Concurrent Verification and Descent}

When multiple developers verify concurrently, their judgment streams may
conflict on shared scopes. The overlap condition
$\overlap{U_i}{U_j}$ for jurisdictions $U_i, U_j$ requires agreement on the
restriction $\res_{U_i \cap U_j}(\J_i) = \res_{U_i \cap U_j}(\J_j)$.

This is precisely the \emph{descent condition} from algebraic geometry.
The descent datum for a covering family $\{U_i \to U\}$ consists of:
\begin{enumerate}
  \item Local objects $s_i \in \J(U_i)$ for each patch.
  \item Gluing isomorphisms $\phi_{ij} : \res_{U_{ij}}(s_i) \xrightarrow{\sim}
        \res_{U_{ij}}(s_j)$ on each double overlap.
  \item The cocycle condition $\phi_{jk} \circ \phi_{ij} = \phi_{ik}$ on triple
        overlaps $U_{ijk} = U_i \cap U_j \cap U_k$.
\end{enumerate}

The descent condition holds iff $\Hcoh{1}(\{U_i\}, \J) = 0$. In the team
context, each developer's judgment stream is a local section, overlap checks
provide gluing isomorphisms, and the cocycle condition ensures transitivity.

% ------------------------------------------------------------------
\section{Jurisdiction Model}\label{sec:jurisdictions}
% ------------------------------------------------------------------

\begin{definition}[Jurisdiction]
A \emph{jurisdiction} for developer $d$ is a sub-site
$\Site_d \hookrightarrow \Site$ such that:
\begin{enumerate}
  \item $\Site_d$ is a full subcategory of $\Site$ on a subset of objects.
  \item The collection $\{\Site_d\}_{d \in D}$ forms a covering family:
        $\bigcup_{d \in D} \Ob(\Site_d) = \Ob(\Site)$.
  \item For each $d$, the developer holds a trust certificate at level
        $\OwnerToken$ for all $U \in \Ob(\Site_d)$.
\end{enumerate}
\end{definition}

The jurisdiction assignment arises naturally from code ownership. In a typical
project with \compSiteTotalFunctions{} functions across \compSiteTotalClasses{}
classes, we partition objects by module ownership. The mean site construction
time of \compSiteMeanBuildMs{}\,ms ensures that jurisdiction boundaries are
computed cheaply.

\begin{definition}[Jurisdiction Overlap]
The \emph{overlap zone} between developers $d_1$ and $d_2$ is:
\[
  \Site_{d_1,d_2} = \Site_{d_1} \times_{\Site} \Site_{d_2}
\]
Objects in $\Site_{d_1,d_2}$ require agreement from both developers for
verification to succeed.
\end{definition}

\subsection{The Jurisdiction Lattice}

The collection of all jurisdictions over a site $\Site$ carries a natural
lattice structure that governs refinement and coarsening of ownership.

\begin{definition}[Jurisdiction Lattice]
Let $D$ be a finite set of developers and $\Site$ a semantic site. The
\emph{jurisdiction lattice} $\JurisCat(\Site, D)$ is the partially ordered set
whose elements are jurisdiction assignments $\alpha : D \to \mathcal{P}(\Ob(\Site))$
satisfying $\bigcup_{d \in D} \alpha(d) = \Ob(\Site)$, ordered by refinement:
$\alpha \leq \beta$ if and only if for every $d \in D$,
$\alpha(d) \subseteq \beta(d)$.
\end{definition}

The lattice has a bottom element $\alpha_{\bot}$, the \emph{minimal partition}
where each object is assigned to exactly one developer, and a top element
$\alpha_{\top}$, the \emph{maximal overlap} where every developer owns every
object. The meet and join operations are:
\begin{align}
  (\alpha \wedge \beta)(d) &= \alpha(d) \cap \beta(d)
  \cup \{U \mid U \notin \textstyle\bigcup_{d'} (\alpha(d') \cap \beta(d'))\},
  \label{eq:juris-meet} \\
  (\alpha \vee \beta)(d) &= \alpha(d) \cup \beta(d).
  \label{eq:juris-join}
\end{align}
The meet requires a redistribution of orphaned objects to maintain the covering
condition; we assign them by proximity in the morphism graph of $\Site$.

\begin{theorem}[Jurisdiction Lattice Properties]\label{thm:juris-lattice}
The jurisdiction lattice $\JurisCat(\Site, D)$ satisfies:
\begin{enumerate}
  \item \textbf{Completeness:} Every finite subset has a meet and join.
  \item \textbf{Covering invariant:} Every element satisfies the covering condition.
  \item \textbf{Overlap monotonicity:} If $\alpha \leq \beta$, then the total
        number of overlap objects under $\alpha$ is at most that under $\beta$.
  \item \textbf{Descent compatibility:} For any $\alpha \in \JurisCat(\Site, D)$,
        the induced covering family $\{\Site_d^{\alpha}\}_{d \in D}$ satisfies
        the descent condition for $\J$ whenever $\Hcoh{1}(\Site, \J) = 0$.
\end{enumerate}
\end{theorem}

\begin{proof}
\textbf{(1)} Completeness follows from finiteness of $\Ob(\Site)$ and $D$.
The join $(\bigvee_k \alpha_k)(d) = \bigcup_k \alpha_k(d)$ preserves coverage.
For the meet, take the pointwise intersection
$\gamma(d) = \bigcap_k \alpha_k(d)$; orphaned objects
$O = \Ob(\Site) \setminus \bigcup_d \gamma(d)$ are redistributed by
shortest-path proximity in the morphism graph.

\textbf{(2)} Maintained by construction: both join (union) and meet (with
redistribution) satisfy $\bigcup_d \alpha(d) = \Ob(\Site)$.

\textbf{(3)} If $\alpha \leq \beta$, then $\alpha(d) \subseteq \beta(d)$,
so $\alpha(d_1) \cap \alpha(d_2) \subseteq \beta(d_1) \cap \beta(d_2)$ and
the overlap count $n_\alpha \leq n_\beta$.

\textbf{(4)} The covering $\{\Site_d^\alpha\}$ refines $\{\Site\}$. By the
comparison theorem for \v{C}ech cohomology, the map
$\Hcoh{1}(\{\Site_d^\alpha\}, \J) \to \Hcoh{1}(\Site, \J) = 0$ is injective,
so the source vanishes. \qed
\end{proof}

\subsection{Jurisdiction Granularity}\label{sec:granularity}

Jurisdictions can be defined at different granularity levels, corresponding to
different choices of objects in $\Site$:

\begin{center}
\begin{tabular}{llll}
\toprule
\textbf{Granularity} & \textbf{Objects} & \textbf{Count} & \textbf{Overlap} \\
\midrule
Module-level   & $\suppCoordModuleCount{}$ modules
               & \suppCoordModulePct{} of coords
               & Low \\
Function-level & $\suppCoordFunctionCount{}$ functions
               & \suppCoordFunctionPct{} of coords
               & Medium \\
Interface-level & $\suppCoordInterfaceCount{}$ interfaces
               & \suppCoordInterfacePct{} of coords
               & High \\
\bottomrule
\end{tabular}
\end{center}

Module-level granularity minimizes overlap but may create bottlenecks.
Function-level enables fine-grained parallelism but increases overlap checks.
Interface-level suits API-driven architectures.

The site has a mean of \compSiteMeanCoords{} coordinates and
\compSiteMeanMorphisms{} morphisms per program, with the morphism distribution
(\compMorphRestriction{} restrictions, \compMorphTransport{} transports,
\compMorphInclusion{} inclusions) determining overlap structure.

\subsection{Jurisdiction Partitioning Algorithm}

We describe how jurisdictions are computed from code ownership metadata.

\begin{algorithm}[t]
\caption{Jurisdiction Partitioning from Code Ownership}
\label{alg:jurisdiction}
\begin{algorithmic}[1]
\REQUIRE Site $\Site$, ownership map $\mathrm{own} : \Ob(\Site) \to 2^D$
\ENSURE Jurisdiction assignment $\alpha : D \to \mathcal{P}(\Ob(\Site))$
\STATE Initialize $\alpha(d) \gets \{U \in \Ob(\Site) \mid d \in \mathrm{own}(U)\}$
        for each $d \in D$
\STATE $O \gets \Ob(\Site) \setminus \bigcup_{d \in D} \alpha(d)$
        \COMMENT{Unowned objects}
\FOR{each $U \in O$}
  \STATE $d^* \gets \arg\min_{d \in D} \min_{V \in \alpha(d)}
          \mathrm{dist}_{\Site}(U, V)$
  \STATE $\alpha(d^*) \gets \alpha(d^*) \cup \{U\}$
\ENDFOR
\STATE Validate: $\bigcup_d \alpha(d) = \Ob(\Site)$
\STATE \textbf{return} $\alpha$
\end{algorithmic}
\end{algorithm}

\begin{example}[Jurisdiction Assignment in \JuGeo]
Consider a web application with modules \texttt{auth}, \texttt{api}, and
\texttt{db}, assigned to developers Alice, Bob, and Carol respectively. The
site $\Site$ has \compSiteTotalFunctions{} function objects. Using the
\JuGeo{} API:

\begin{lstlisting}[language=Python,basicstyle=\small\ttfamily,frame=single,
  xleftmargin=1em]
from jugeo import SiteBuilder, DescentEngine, DescentConfiguration
from jugeo import DescentStrategy

source_code = load_project("webapp")
site = SiteBuilder(source_code).build()

# Define jurisdictions from CODEOWNERS
ownership = {
    "auth/*": "alice", "api/*": "bob", "db/*": "carol"
}
jurisdictions = site.partition_by_ownership(ownership)

# Each developer verifies their jurisdiction
for dev, sub_site in jurisdictions.items():
    cover = sub_site.topology.covers[0]
    config = DescentConfiguration(
        strategy=DescentStrategy.EAGER, depth_limit=5
    )
    engine = DescentEngine(configuration=config)
    result = engine.run(cover, sub_site.local_sections())
    assert result.success  # Local verification passes

# Merge across jurisdictions
merged = site.merge_jurisdictions(jurisdictions)
assert merged.obstruction_rank == 0  # Global gluing succeeds
\end{lstlisting}
\end{example}

\subsection{Covering Conditions}

We require the jurisdiction cover to satisfy a \v{C}ech-type condition. The
\v{C}ech complex $\Cech(\{U_d\})$ has:
\begin{itemize}
  \item $C^0$: judgments on individual jurisdictions,
  \item $C^1$: agreement constraints on pairwise overlaps,
  \item $C^2$: triple-overlap coherence (transitivity of agreement).
\end{itemize}
The verification succeeds globally when $\Hcoh{1}(\Site, \mathcal{J}) = 0$,
which our experiments confirm: all \expHOneZero{} programs achieve vanishing
first cohomology.

More precisely, the \v{C}ech complex in degree $p$ for a team of developers
$D = \{d_0, \ldots, d_n\}$ is:
\[
  \Cech^p(\{U_d\}_{d \in D}, \J) =
  \prod_{\substack{d_0 < d_1 < \cdots < d_p \\ d_i \in D}}
  \J\!\left(\bigcap_{i=0}^p U_{d_i}\right)
\]
The coboundary $\delta^0 : \Cech^0 \to \Cech^1$ sends a family of local
judgments $(s_d)_{d \in D}$ to the family of ``disagreements''
$(\delta^0 s)_{d_1, d_2} = \res_{U_{d_1 d_2}}(s_{d_2}) - \res_{U_{d_1 d_2}}(s_{d_1})$
where the subtraction is in the abelianized trust algebra. The kernel of
$\delta^0$ is precisely $\Hcoh{0}$, the global sections, and the cokernel
$\ker \delta^1 / \mathrm{im}\,\delta^0$ is $\Hcoh{1}$, the obstruction group.

% ------------------------------------------------------------------
\section{Authority Delegation Protocol}\label{sec:delegation}
% ------------------------------------------------------------------

\subsection{The Delegation Category}

We formalize authority delegation as a category, providing a compositional
framework for reasoning about trust flow in teams.

\begin{definition}[Delegation Category]
The \emph{delegation category} $\DelegCat(\Site, D)$ has:
\begin{itemize}
  \item \textbf{Objects:} Pairs $(d, U)$ with $d \in D$ and
        $U \in \Ob(\Site_d)$.
  \item \textbf{Morphisms:} $\delta : (d_1, U) \to (d_2, V)$ when $d_1$
        delegates over $U$ to $d_2$ for $V \subseteq U$, with trust
        $\tatten(\trust_{d_1}|_U)$.
  \item \textbf{Composition:} $\delta_2 \circ \delta_1 : (d_1, U) \to (d_3, W)$
        carries trust $\tatten^2(\trust_{d_1}|_U)$.
  \item \textbf{Identity:} $\id_{(d,U)}$ is self-delegation with no attenuation.
\end{itemize}
\end{definition}

\begin{proposition}[Well-Definedness]
$\DelegCat(\Site, D)$ is a category: composition is associative (since
$\tatten$ composes associatively and scope inclusion is transitive) and
identities are neutral.
\end{proposition}

\subsection{Delegation Morphisms}

A delegation from developer $d_1$ to $d_2$ over scope $U$ is a trust morphism:
\[
  \delta_{d_1 \to d_2}^U : \trust_{d_1}|_U \xrightarrow{\tatten}
  \trust_{d_2}|_U
\]
The attenuation $\tatten$ ensures the delegatee operates at a trust level
strictly below the delegator. If $d_1$ holds $\Tsolver$ trust on $U$, then
$d_2$ receives at most $\Tcopilot$ trust via delegation.

\begin{algorithm}[t]
\caption{Authority Delegation via \texttt{public\_alignment}}
\label{alg:delegation}
\begin{algorithmic}[1]
\REQUIRE Developer $d_1$ (delegator), $d_2$ (delegatee), scope $U$
\ENSURE Delegation token $\delta$ with attenuated trust
\STATE $\trust_1 \gets$ \texttt{public\_alignment}$(d_1, U)$
\STATE $\trust_2 \gets \tatten(\trust_1)$ \COMMENT{Attenuate one level}
\STATE $\delta \gets (\text{from}: d_1, \text{to}: d_2, \text{scope}: U,
       \text{trust}: \trust_2, \text{ttl}: T)$
\STATE Register $\delta$ in $\AuthGraph$
\STATE \textbf{return} $\delta$
\end{algorithmic}
\end{algorithm}

\subsection{Delegation Chains and Composition}

Delegations compose: if $d_1$ delegates to $d_2$ and $d_2$ delegates to $d_3$,
the composed delegation attenuates twice. We bound chain length at 3 to prevent
trust from dropping below $\Tunver$.

\begin{theorem}[Delegation Composition Soundness]\label{thm:deleg-compose}
Let $\delta_1 : (d_1, U_1) \to (d_2, U_2)$ and
$\delta_2 : (d_2, U_2) \to (d_3, U_3)$ be delegation morphisms in
$\DelegCat(\Site, D)$ with $U_3 \subseteq U_2 \subseteq U_1$. Then:
\begin{enumerate}
  \item The composed delegation $\delta_2 \circ \delta_1$ has trust level
        $\trust_{d_3} = \tatten^2(\trust_{d_1}|_{U_1})$.
  \item If $\trust_{d_1} = \Tsolver$, then
        $\trust_{d_3} \in \{\Tunver, \Tcopilot\}$.
  \item For any chain $d_1 \to d_2 \to \cdots \to d_k$ with $k \leq 3$,
        $\trust_{d_k} \neq \Tcontra$.
  \item The delegation functor $\DelegCat(\Site, D) \to \Talg^{\op}$ sending
        each delegation to its trust level is a contravariant functor.
\end{enumerate}
\end{theorem}

\begin{proof}
\textbf{(1)} By definition, the trust level of
$\delta_2 \circ \delta_1$ is $\tatten(\tatten(\trust_{d_1}|_{U_1})) =
\tatten^2(\trust_{d_1}|_{U_1})$.

\textbf{(2)} The attenuation map on the trust lattice
$\Tproof \xrightarrow{\tatten} \Tsolver \xrightarrow{\tatten} \Truntime
\xrightarrow{\tatten} \Toracle \xrightarrow{\tatten} \Tcopilot
\xrightarrow{\tatten} \Tunver$ drops one level per application. Starting from
$\Tsolver$, two attenuations yield $\Tcopilot$ or $\Tunver$.

\textbf{(3)} By induction on chain length $k$. For $k = 1$,
$\tatten(\trust) \neq \Tcontra$ since $\tatten$ maps
$\{\Tunver, \ldots, \Tproof\}$ to $\{\Tunver, \ldots, \Tsolver\}$.
For $k + 1 \leq 3$, the inductive hypothesis gives
$\trust_{d_k} \geq \Tunver$, and $\tatten(\Tunver) = \Tunver \neq \Tcontra$.

\textbf{(4)} The functor sends $\delta : (d_1, U) \to (d_2, V)$ to
$\tatten(\trust_{d_1}) \tleq \trust_{d_1}$ in $\Talg^{\op}$, preserving
composition by order-preservation of $\tatten$. \qed
\end{proof}

\begin{proposition}[Delegation Soundness]
For any delegation chain $d_1 \to d_2 \to \cdots \to d_k$ with $k \leq 3$,
the resulting trust level satisfies $\trust_{d_k} \tleq \trust_{d_1}$ and
$\trust_{d_k} \neq \Tcontra$.
\end{proposition}

\subsection{Delegation Patterns}\label{sec:deleg-patterns}

We identify three common delegation patterns that arise in team workflows:

\paragraph{Peer Delegation.}
Two developers at the same trust level delegate to each other for cross-review.
If both hold $\Tsolver$ trust, the delegation creates symmetric morphisms
$\delta_{d_1 \to d_2}^U$ and $\delta_{d_2 \to d_1}^V$ with both delegatees
operating at $\Tcopilot$. This pattern supports pair programming and code
review.

\paragraph{Hierarchical Delegation.}
A team lead with $\Tsolver$ trust delegates to juniors at $\Tcopilot$, who may
delegate to interns at $\Tunver$ (chain length 2). This three-tier model bounds
delegation chains at 3 levels.

\paragraph{Scope-Restricted Delegation.}
A developer delegates over a strict sub-scope: $d_1$ owns \texttt{auth} but
delegates only \texttt{auth.oauth} to $d_2$, retaining full authority elsewhere.

\begin{example}[Delegation in the \JuGeo{} API]
The following code demonstrates hierarchical delegation with revocation:

\begin{lstlisting}[language=Python,basicstyle=\small\ttfamily,frame=single,
  xleftmargin=1em]
from jugeo import SiteBuilder, DelegationRegistry, TrustLevel

site = SiteBuilder(source_code).build()
registry = DelegationRegistry(site)

# Team lead Alice delegates auth module to Bob
token_ab = registry.delegate(
    delegator="alice", delegatee="bob",
    scope=site.scope("auth"),
    ttl=3600  # 1 hour expiration
)
assert token_ab.trust_level == TrustLevel.COPILOT

# Bob sub-delegates oauth to Carol (chain length 2)
token_bc = registry.delegate(
    delegator="bob", delegatee="carol",
    scope=site.scope("auth.oauth"),
    parent_token=token_ab
)
assert token_bc.trust_level == TrustLevel.UNVERIFIED

# Revoking Alice->Bob cascades to Bob->Carol
registry.revoke(token_ab)
assert not registry.is_active(token_bc)
\end{lstlisting}
\end{example}

\subsection{Revocation}

Delegation tokens carry a time-to-live (TTL) and can be explicitly revoked.
Revocation propagates through the $\AuthGraph$: revoking $\delta_{d_1 \to d_2}$
automatically invalidates all downstream delegations from $d_2$.

The revocation semantics are formalized as follows. The $\AuthGraph$ is a
directed acyclic graph (DAG) whose vertices are active delegation tokens and
whose edges represent parent-child relationships. Revoking a token $\delta$
triggers a depth-first traversal of the sub-DAG rooted at $\delta$, marking
each descendant as revoked. Judgments produced under revoked tokens are demoted
to $\Tunver$ trust until re-verified.

\begin{proposition}[Revocation Completeness]
After revoking delegation token $\delta_0$, for every descendant
$\delta \in \AuthGraph$ and judgment $\J$ produced under $\delta$,
$\trust(\J) \tleq \Tunver$.
\end{proposition}

% ------------------------------------------------------------------
\section{Merge Semantics for Concurrent Judgments}\label{sec:merge}
% ------------------------------------------------------------------

\subsection{The Merge Functor}

We formalize the merge operation as a functor from the category of local
judgment families to the category of global judgments.

\begin{definition}[Merge Functor]
Let $\mathcal{U} = \{\Site_d\}_{d \in D}$ be a jurisdiction cover. The
\emph{merge functor} $\MergeFun : \prod_{d \in D} \Sh(\Site_d) \to \Sh(\Site)$
sends a family of local sheaves $(\J_d)_{d \in D}$ with compatible restrictions
on overlaps to the unique global sheaf $\J$ satisfying $\J|_{\Site_d} = \J_d$.
The functor is well-defined whenever the \v{C}ech cohomology
$\Hcoh{1}(\mathcal{U}, \J) = 0$.
\end{definition}

The merge functor $\MergeFun$ is dual to the restriction
$\res^* : \Sh(\Site) \to \prod_d \Sh(\Site_d)$: $\MergeFun \circ \res^* = \id$
on $\Sh(\Site)$, and $\res^* \circ \MergeFun = \id$ on compatible families.

\subsection{Judgment Streams}

Each developer $d$ produces a stream of judgments
$\mathcal{J}_d = \{\J_{d,1}, \J_{d,2}, \ldots\}$ over their jurisdiction.
When developers work concurrently, the system must merge these streams into a
consistent global judgment.

\begin{definition}[Judgment Merge]
The merge of judgments $\J_1, \J_2$ on the same coordinate $\coord$ is:
\[
  \MergeOp(\J_1, \J_2) = \Jud{\coord}
    {\prop_1 \wedge \prop_2}
    {\carrier_1 \cup \carrier_2}
    {\evid_1 \tjoin \evid_2}
    {\oblig_1 \cup \oblig_2}
    {\obstr_1 \cup \obstr_2}
    {\trust_1 \tjoin \trust_2}
    {\prov_1 \| \prov_2}
\]
where $\|$ denotes provenance concatenation.
\end{definition}

\begin{theorem}[Merge Associativity and Commutativity]\label{thm:merge-props}
The merge operation $\MergeOp$ satisfies:
\begin{enumerate}
  \item \textbf{Commutativity:} $\MergeOp(\J_1, \J_2) = \MergeOp(\J_2, \J_1)$.
  \item \textbf{Associativity:}
        $\MergeOp(\MergeOp(\J_1, \J_2), \J_3) =
         \MergeOp(\J_1, \MergeOp(\J_2, \J_3))$.
  \item \textbf{Idempotence:} $\MergeOp(\J, \J) = \J$.
  \item \textbf{Trust monotonicity:}
        $\trust(\MergeOp(\J_1, \J_2)) \geq \min(\trust(\J_1), \trust(\J_2))$
        under $\tleq$.
\end{enumerate}
\end{theorem}

\begin{proof}
\textbf{(1)} Commutativity follows from commutativity of each component:
$\wedge$ on propositions, $\cup$ on sets, and
$\tjoin$ on evidence and trust (by the axioms of $\Talg$).

\textbf{(2)} Associativity follows component-wise from associativity of
$\wedge$, $\cup$, and $\tjoin$ (the semilattice axiom in $\Talg$).

\textbf{(3)} Idempotence: $\prop \wedge \prop = \prop$,
$\carrier \cup \carrier = \carrier$,
$\evid \tjoin \evid = \evid$ (by idempotence of $\tjoin$).

\textbf{(4)} Trust monotonicity follows from $\tjoin$ being a join:
$\trust_1 \tjoin \trust_2 \geq \trust_1$ and
$\trust_1 \tjoin \trust_2 \geq \trust_2$ in $(\Talg, \tleq)$. \qed
\end{proof}

\subsection{Conflict Detection}

A conflict arises when $\J_1$ and $\J_2$ on the same $\coord$ produce
contradictory propositions. Formally, a conflict exists when
$\prop_1 \wedge \prop_2$ is unsatisfiable in the SMT theory $\QFLIA + \QFUF$.

\begin{definition}[Merge Conflict]
A \emph{merge conflict} at coordinate $\coord$ between judgments $\J_1, \J_2$
is a tuple $(\coord, \prop_1, \prop_2, \sigma)$ where $\sigma$ is a certificate
of unsatisfiability for $\prop_1 \wedge \prop_2$. The \emph{conflict class}
$[\sigma] \in \Hcoh{1}(\{\Site_{d_1}, \Site_{d_2}\}, \J)$ is the cohomology
class representing this obstruction.
\end{definition}

\begin{algorithm}[t]
\caption{Conflict Resolution via \texttt{orchestrate\_verification}}
\label{alg:conflict}
\begin{algorithmic}[1]
\REQUIRE Judgment streams $\mathcal{J}_{d_1}, \mathcal{J}_{d_2}$
\ENSURE Merged stream $\mathcal{J}_{\text{merged}}$ or obstruction report
\FOR{each shared coordinate $\coord \in \Site_{d_1,d_2}$}
  \STATE $\J_1 \gets \mathcal{J}_{d_1}(\coord)$,
         $\J_2 \gets \mathcal{J}_{d_2}(\coord)$
  \STATE $\phi \gets$ \texttt{orchestrate\_verification}$(\prop_1 \wedge \prop_2)$
  \IF{$\phi$ is SAT}
    \STATE $\mathcal{J}_{\text{merged}}(\coord) \gets \MergeOp(\J_1, \J_2)$
  \ELSE
    \STATE Record obstruction at $\coord$; attempt \REPAIR
  \ENDIF
\ENDFOR
\end{algorithmic}
\end{algorithm}

\subsection{$n$-Way Merge}\label{sec:nway-merge}

For triple overlaps $\Site_{d_1,d_2,d_3}$, we extend the merge to a three-way
operation. The cocycle condition ensures that
$\MergeOp(\MergeOp(\J_1, \J_2), \J_3) = \MergeOp(\J_1, \MergeOp(\J_2, \J_3))$,
which holds because $\tjoin$ is associative and commutative.

More generally, we define the $n$-way merge for an arbitrary number of
developers with overlapping jurisdictions.

\begin{algorithm}[t]
\caption{$n$-Way Judgment Merge}
\label{alg:nway-merge}
\begin{algorithmic}[1]
\REQUIRE Jurisdiction cover $\{(\Site_{d_i}, \J_{d_i})\}_{i=1}^n$
\ENSURE Global merged judgment $\J_{\mathrm{global}}$ or obstruction report
\STATE Compute the nerve $\nerve(\mathcal{U})$ of the cover
\FOR{each $p$-simplex $\sigma = [d_{i_0}, \ldots, d_{i_p}]$ in $\nerve(\mathcal{U})$}
  \STATE $U_\sigma \gets \bigcap_{k=0}^p \Site_{d_{i_k}}$
  \IF{$U_\sigma \neq \emptyset$ and $p = 1$}
    \STATE Check pairwise compatibility:
           $\res_{U_\sigma}(\J_{d_{i_0}}) \stackrel{?}{=}
            \res_{U_\sigma}(\J_{d_{i_1}})$
    \IF{incompatible}
      \STATE $\phi \gets$ \texttt{encode\_for\_solver}$(\prop_{i_0} \wedge \prop_{i_1})$
      \IF{$\phi$ is UNSAT}
        \STATE Record 1-cocycle obstruction at $\sigma$
      \ENDIF
    \ENDIF
  \ENDIF
  \IF{$U_\sigma \neq \emptyset$ and $p = 2$}
    \STATE Check cocycle condition on triple overlap
  \ENDIF
\ENDFOR
\IF{no obstructions recorded}
  \STATE Apply sheaf gluing to construct $\J_{\mathrm{global}}$
  \STATE \textbf{return} $\J_{\mathrm{global}}$
\ELSE
  \STATE \textbf{return} obstruction report with $\Hcoh{1}$ classes
\ENDIF
\end{algorithmic}
\end{algorithm}

The $n$-way merge algorithm traverses the nerve of the jurisdiction cover,
checking compatibility on each simplex. For our benchmarks with 2--4
jurisdictions, the nerve has at most 6 edges and 4 triangles.

\subsection{Connection to \v{C}ech Cohomology}

The merge process has a precise cohomological interpretation. Given the
jurisdiction cover $\mathcal{U} = \{U_d\}_{d \in D}$ and a family of local
judgments $(s_d \in \J(U_d))_{d \in D}$, define the \v{C}ech 1-cochain:
\[
  c^1_{d_1, d_2} = \res_{U_{d_1 d_2}}(s_{d_2}) \tjoin
  (\tatten \circ \res_{U_{d_1 d_2}}(s_{d_1}))^{-1}
\]
where the ``inverse'' is taken in the abelianized trust group. The family
$(s_d)$ glues to a global section if and only if $c^1$ is a coboundary, \ie
$[c^1] = 0 \in \Hcoh{1}(\mathcal{U}, \J)$.

\begin{example}[Merge in the \JuGeo{} API]
The following code demonstrates the $n$-way merge workflow:

\begin{lstlisting}[language=Python,basicstyle=\small\ttfamily,frame=single,
  xleftmargin=1em]
from jugeo import SiteBuilder, MergeEngine, DescentEngine
from jugeo import DescentConfiguration, DescentStrategy

site = SiteBuilder(source_code).build()
jurisdictions = site.partition_by_ownership(ownership_map)

# Each developer runs descent independently
local_results = {}
for dev, sub_site in jurisdictions.items():
    config = DescentConfiguration(
        strategy=DescentStrategy.EAGER, depth_limit=5
    )
    engine = DescentEngine(configuration=config)
    local_results[dev] = engine.run(
        sub_site.topology.covers[0],
        sub_site.local_sections()
    )

# n-way merge across all jurisdictions
merger = MergeEngine(site)
global_result = merger.merge_all(local_results)

if global_result.success:
    print(f"Global merge succeeded, "
          f"obstruction rank: {global_result.obstruction_rank}")
    # obstruction_rank == 0 means H^1 = 0
else:
    for obs in global_result.obstructions:
        print(f"Conflict at {obs.coordinate}: "
              f"{obs.cohomology_class}")
\end{lstlisting}
\end{example}

% ------------------------------------------------------------------
\section{Cohomological Analysis of Team Conflicts}\label{sec:cohomology}
% ------------------------------------------------------------------

We now develop the cohomological framework for analyzing team conflicts in
full generality. The key idea is that the team \v{C}ech complex encodes all
possible obstructions to merging concurrent verification results, and the
vanishing of its cohomology guarantees successful global merge.

\subsection{The Team \v{C}ech Complex}

\begin{definition}[Team \v{C}ech Complex]
Let $D = \{d_1, \ldots, d_n\}$ be a development team and
$\mathcal{U} = \{\Site_{d_i}\}_{i=1}^n$ the jurisdiction cover.
The \emph{team \v{C}ech complex} $\TeamCech^\bullet(\mathcal{U}, \J)$ is the
cochain complex:
\[
  0 \to \Cech^0(\mathcal{U}, \J) \xrightarrow{\delta^0}
  \Cech^1(\mathcal{U}, \J) \xrightarrow{\delta^1}
  \Cech^2(\mathcal{U}, \J) \xrightarrow{\delta^2} \cdots
\]
where:
\begin{align*}
  \Cech^0(\mathcal{U}, \J) &=
    \prod_{d \in D} \J(\Site_d)
    & &\text{(local judgments per developer),} \\
  \Cech^1(\mathcal{U}, \J) &=
    \prod_{d_i < d_j} \J(\Site_{d_i} \cap \Site_{d_j})
    & &\text{(pairwise overlap data),} \\
  \Cech^2(\mathcal{U}, \J) &=
    \prod_{d_i < d_j < d_k} \J(\Site_{d_i} \cap \Site_{d_j} \cap \Site_{d_k})
    & &\text{(triple coherence).}
\end{align*}
\end{definition}

The coboundary maps are defined by the alternating restriction formula.
For a 0-cochain $(s_d)_{d \in D} \in \Cech^0$:
\[
  (\delta^0 s)_{d_i, d_j} =
  \res_{\Site_{d_i} \cap \Site_{d_j}}(s_{d_j})
  \tjoin \bigl(\res_{\Site_{d_i} \cap \Site_{d_j}}(s_{d_i})\bigr)^{\tatten}
\]
where the superscript $\tatten$ denotes the ``trust difference'' in the
abelianized trust group $\Talg / [\Talg, \Talg]$. When the trust algebra is
commutative (as in \JuGeo), this simplifies to:
\[
  (\delta^0 s)_{d_i, d_j} =
  \res_{\Site_{d_i d_j}}(s_{d_j}) \tjoin
  \tatten\bigl(\res_{\Site_{d_i d_j}}(s_{d_i})\bigr).
\]

For a 1-cochain $(c_{d_i, d_j}) \in \Cech^1$:
\[
  (\delta^1 c)_{d_i, d_j, d_k} =
  \res_{\Site_{d_i d_j d_k}}(c_{d_j, d_k})
  \tjoin \tatten\bigl(\res_{\Site_{d_i d_j d_k}}(c_{d_i, d_k})\bigr)
  \tjoin \res_{\Site_{d_i d_j d_k}}(c_{d_i, d_j}).
\]

\subsection{Computation of $\Hcoh{0}$ and $\Hcoh{1}$}

\begin{theorem}[$\Hcoh{0}$ as Global Agreements]\label{thm:h0}
The zeroth \v{C}ech cohomology
$\Hcoh{0}(\mathcal{U}, \J) = \ker \delta^0$ is isomorphic to
$\J(\Site)$, the space of global judgments. That is, $\Hcoh{0}$ classifies
exactly those families of local judgments that are globally consistent.
\end{theorem}

\begin{proof}
An element $(s_d)_{d \in D} \in \ker \delta^0$ satisfies
$\res_{\Site_{d_i d_j}}(s_{d_j}) = \res_{\Site_{d_i d_j}}(s_{d_i})$ for all
pairs $d_i, d_j$ (since $\tjoin$ with the attenuated inverse yields the neutral
element only when the restrictions agree). This is precisely the gluing
condition for the sheaf $\J$. By the sheaf axiom, such a compatible family
extends to a unique global section $s \in \J(\Site)$. Conversely, any global
section restricts to a compatible family. The map
$\J(\Site) \to \ker \delta^0$ given by $s \mapsto (\res_{\Site_d}(s))_{d \in D}$
is therefore a bijection. \qed
\end{proof}

\begin{theorem}[$\Hcoh{1}$ as Merge Obstructions]\label{thm:h1}
The first \v{C}ech cohomology
$\Hcoh{1}(\mathcal{U}, \J) = \ker \delta^1 / \mathrm{im}\,\delta^0$ classifies
equivalence classes of merge obstructions. Specifically:
\begin{enumerate}
  \item Each nonzero class $[\sigma] \in \Hcoh{1}$ corresponds to a family of
        local judgments that are pairwise compatible on overlaps (satisfying
        $\delta^1 \sigma = 0$, the cocycle condition) but cannot be extended to
        a global judgment.
  \item $\Hcoh{1}(\mathcal{U}, \J) = 0$ if and only if the merge functor
        $\MergeFun$ is well-defined on all compatible families.
  \item Two obstructions $\sigma, \sigma'$ are equivalent ($[\sigma] = [\sigma']$)
        if they differ by a coboundary, \ie a ``gauge transformation'' arising from
        modifying local judgments within their equivalence classes.
\end{enumerate}
\end{theorem}

\begin{proof}
\textbf{(1)} A 1-cocycle $\sigma \in \ker \delta^1$ satisfies the cocycle
condition on triple overlaps. If $\sigma = \delta^0 s$, the family glues
globally. Otherwise, $[\sigma] \neq 0$ represents a genuine obstruction.

\textbf{(2)} $\MergeFun$ is well-defined on all compatible families iff every
cocycle is a coboundary, \ie $\Hcoh{1} = 0$.

\textbf{(3)} Two cocycles $\sigma, \sigma'$ differing by $\delta^0 t$
represent the same obstruction class since $t$ is a local adjustment. \qed
\end{proof}

\subsection{Spectral Sequence for Hierarchical Teams}

In hierarchical team structures, we refine the analysis using a spectral
sequence.

\begin{construction}[Leray Spectral Sequence for Team Hierarchy]
Let $\pi : D \to L$ be the team hierarchy map sending each developer $d$ to
their team lead $\ell = \pi(d)$. This induces a factorization of the
jurisdiction cover: the \emph{coarse cover} $\{V_\ell\}_{\ell \in L}$ where
$V_\ell = \bigcup_{\pi(d) = \ell} \Site_d$ (each lead's combined jurisdiction),
and for each lead $\ell$, the \emph{fine cover}
$\{U_d\}_{\pi(d) = \ell}$ of $V_\ell$.

The Leray spectral sequence for this two-stage cover has:
\[
  E_2^{p,q} = \Hcoh{p}(\{V_\ell\}_{\ell \in L},\;
  \mathcal{H}^q(\{U_d\}_d, \J))
  \implies
  \Hcoh{p+q}(\mathcal{U}, \J)
\]
where $\mathcal{H}^q$ is the presheaf of local cohomology groups.
\end{construction}

\begin{corollary}[Hierarchical Vanishing]
If $\Hcoh{1}(\{U_d\}_{\pi(d)=\ell}, \J|_{V_\ell}) = 0$ for each team lead
$\ell$ (each sub-team merges successfully) and
$\Hcoh{1}(\{V_\ell\}_{\ell \in L}, \J) = 0$ (the leads' combined
jurisdictions merge successfully), then $\Hcoh{1}(\mathcal{U}, \J) = 0$
(the full team merge succeeds).
\end{corollary}

\begin{proof}
Under the hypotheses, $E_2^{1,0} = \Hcoh{1}(\{V_\ell\}, \Hcoh{0}(\{U_d\}, \J)) =
\Hcoh{1}(\{V_\ell\}, \J) = 0$ and
$E_2^{0,1} = \Hcoh{0}(\{V_\ell\}, \mathcal{H}^1) = 0$ since each
$\mathcal{H}^1(\{U_d\}_{\pi(d)=\ell}, \J) = 0$. The spectral sequence
degenerates to give $\Hcoh{1}(\mathcal{U}, \J) = 0$. \qed
\end{proof}

This result has practical significance: hierarchical teams can
verify in two stages---first each sub-team merges internally, then leads merge
across sub-teams---with guaranteed global soundness.

\subsection{Long Exact Sequence for Sub-Teams}\label{sec:long-exact}

We also establish a long exact sequence relating team and sub-team cohomology.
Let $D' \subset D$ be a sub-team and $D'' = D \setminus D'$.

\begin{proposition}[Mayer--Vietoris for Team Cohomology]\label{prop:mayer-vietoris}
Let $V' = \bigcup_{d \in D'} \Site_d$ and $V'' = \bigcup_{d \in D''} \Site_d$.
There is a long exact sequence:
\[
  0 \to \Hcoh{0}(\Site, \J) \to
  \Hcoh{0}(V', \J) \oplus \Hcoh{0}(V'', \J) \to
  \Hcoh{0}(V' \cap V'', \J) \xrightarrow{\;\partial\;}
  \Hcoh{1}(\Site, \J) \to \cdots
\]
The connecting homomorphism $\partial$ sends a global judgment on the overlap
$V' \cap V''$ that extends over $V'$ and $V''$ separately but not globally to
its obstruction class in $\Hcoh{1}(\Site, \J)$.
\end{proposition}

This exact sequence is the Mayer--Vietoris sequence for the open cover
$\{V', V''\}$ of $\Site$. When the sub-teams agree on all overlap judgments
(\ie the restriction map to $\Hcoh{0}(V' \cap V'', \J)$ is surjective), then
$\partial = 0$ and $\Hcoh{1}(\Site, \J)$ receives no contribution from the
team split. In practice, well-defined module boundaries lead to
obstruction-free team merges.

% ------------------------------------------------------------------
\section{Workflow Orchestration}\label{sec:orchestration}
% ------------------------------------------------------------------

\subsection{Pull Request Verification}

When a developer opens a pull request, the \TeamSite{} computes the
\emph{diff jurisdiction}---the subset of $\Site$ affected by the changes.
The system runs \texttt{discovery\_pipeline} on the diff jurisdiction and
verifies only the affected coordinates, reusing cached judgments for
unchanged regions.

The mean orchestration time of \compProfOrchestrateVerificationMeanMs{}\,ms per
call ensures that PR verification fits within interactive feedback loops. With
\compTrustSolverdischarged{} propositions discharged by the solver at a rate of
\compTrustSolverdischargedPct{}, the system provides high-confidence results.

\subsection{Branch-Parallel Verification}

Feature branches define natural jurisdictions. Each branch's verification runs
independently, with the merge point triggering conflict detection on the
overlap. The process follows:

\begin{enumerate}[leftmargin=*]
  \item \textbf{Fork}: Developer creates branch; jurisdiction inherits from
        parent.
  \item \textbf{Verify}: \texttt{orchestrate\_verification} runs on modified
        scopes using the eager descent strategy (success rate:
        \compDescentEagerRate{}).
  \item \textbf{Merge}: \MergeEngine{} computes overlap, checks cocycle
        condition, and produces merged judgments.
  \item \textbf{Certify}: Merged judgments are certified at $\Tsolver$ level
        via \texttt{encode\_for\_solver}.
\end{enumerate}

\subsection{Role-Based Access}

We define three roles with decreasing authority:
\begin{center}
\begin{tabular}{lll}
\toprule
\textbf{Role} & \textbf{Trust Level} & \textbf{Capabilities} \\
\midrule
$\OwnerToken$ & $\Tsolver$ & Full verification, delegation \\
$\ReviewToken$ & $\Tcopilot$ & Read judgments, challenge claims \\
$\DelegToken$ & $\tatten(\trust_{\text{delegator}})$ & Scoped verification \\
\bottomrule
\end{tabular}
\end{center}

% ------------------------------------------------------------------
\section{Soundness of Team Verification}\label{sec:soundness}
% ------------------------------------------------------------------

We now prove the main soundness theorems for team verification, connecting the
categorical and cohomological machinery to concrete verification guarantees.

\begin{theorem}[Team Merge Soundness]\label{thm:team-soundness}
Let $\{\Site_d\}_{d \in D}$ be a jurisdiction cover of $\Site$ with pairwise
overlaps $\Site_{d_i, d_j}$. If each developer $d$ produces a verified
judgment sheaf $\mathcal{J}_d$ on $\Site_d$, and the merge
$\MergeOp(\J_{d_i}, \J_{d_j})$ succeeds on all overlaps, then the global
merged sheaf $\mathcal{J}$ is a valid judgment sheaf on $\Site$ with
$\Hcoh{1}(\Site, \mathcal{J}) = 0$.
\end{theorem}

\begin{proof}
\textbf{Step 1: Local sections.}
Each $\mathcal{J}_d$ is a section of the judgment sheaf over $\Site_d$.
By hypothesis, $s_d \in \J(\Site_d)$ with $\trust(s_d) \geq \Tcopilot$ for all $d \in D$.

\textbf{Step 2: Cocycle condition.}
The merge condition on overlaps ensures that for all pairs $d_i, d_j \in D$,
\[
  \res_{\Site_{d_i d_j}}(s_{d_i}) \tjoin \res_{\Site_{d_i d_j}}(s_{d_j})
  \neq \Tcontra.
\]
The 1-cochain $c_{d_i, d_j} = \res_{\Site_{d_i d_j}}(s_{d_j}) \tjoin
\tatten(\res_{\Site_{d_i d_j}}(s_{d_i}))$ is a coboundary $\delta^0 s$,
hence $[c] = 0 \in \Hcoh{1}(\mathcal{U}, \J)$. The cocycle condition on
triple overlaps $\delta^1 c = 0$ follows from associativity of $\tjoin$
(Theorem~\ref{thm:merge-props}).

\textbf{Step 3: Gluing.}
By the sheaf gluing axiom, local sections satisfying the cocycle
condition extend to a unique global section $s \in \J(\Site)$ with trust
$\bigwedge_{d \in D} \trust(s_d) \geq \Tcopilot$. The vanishing
$\Hcoh{1}(\Site, \J) = 0$ follows from exactness of the
\v{C}ech complex at degree 1. \qed
\end{proof}

\begin{corollary}[Delegation Preserves Soundness]
Authority delegation does not introduce unsound judgments: any judgment produced
under delegation at trust $\trust'$ satisfies $\trust' \tleq \trust_{\text{owner}}$.
\end{corollary}

\begin{proof}
By Theorem~\ref{thm:deleg-compose}, delegation morphisms in $\DelegCat$
strictly attenuate trust. A judgment $\J$ produced under delegation
$\delta_{d_1 \to d_2}^U$ carries trust
$\trust(\J) = \tatten(\trust_{d_1}|_U) \tleq \trust_{d_1}|_U$.
Since the merge operation $\MergeOp$ uses the conservative join $\tjoin$,
the merged global judgment has trust
$\trust(\MergeOp(\J_1, \J_2)) = \trust(\J_1) \tjoin \trust(\J_2) \geq
\min(\trust(\J_1), \trust(\J_2))$,
preserving the trust bound. \qed
\end{proof}

\begin{corollary}[Hierarchical Team Soundness]
For a hierarchical team with leads $L$ and sub-teams
$\{D_\ell\}_{\ell \in L}$, if each sub-team's internal merge succeeds and
the leads' cross-team merge succeeds, the full merge is sound with
$\Hcoh{1}(\Site, \J) = 0$.
\end{corollary}

\begin{proof}
Immediate from the hierarchical vanishing corollary
(Section~\ref{sec:cohomology}) and Theorem~\ref{thm:team-soundness}. \qed
\end{proof}

% ------------------------------------------------------------------
\section{Lean 4 Formalization}\label{sec:lean}
% ------------------------------------------------------------------

We formalize the key properties in \LEAN, providing machine-checked proofs of
jurisdiction lattice properties, delegation soundness, and merge correctness.

\begin{lstlisting}[basicstyle=\small\ttfamily,
  keywordstyle=\bfseries,frame=single,xleftmargin=1em,
  morekeywords={structure,theorem,def,where,forall,exists,by,apply,exact,
  induction,with,unfold,simp,Type,Set,Fintype,Finset}]
-- Jurisdiction as a sub-site
structure Jurisdiction (S : Site) where
  objects : Set S.Obj
  closed_under_morph : forall {U V : S.Obj},
    U \in objects -> S.Mor U V -> V \in objects

-- Jurisdiction cover
structure JurisdictionCover (S : Site) where
  devs : Type
  assign : devs -> Jurisdiction S
  covers : forall (U : S.Obj),
    exists d, U \in (assign d).objects

-- Delegation with attenuation
structure Delegation (T : TrustAlgebra) where
  from : Developer
  to : Developer
  scope : Set Site.Obj
  trust_level : T.Level
  attenuated : trust_level <=_T T.owner_level from

-- Delegation category
structure DelegationCategory (S : Site) (T : TrustAlgebra) where
  objects : Type := Developer x Set S.Obj
  morphisms : objects -> objects -> Type
  compose : {a b c : objects} ->
    morphisms a b -> morphisms b c -> morphisms a c
  id : (a : objects) -> morphisms a a

-- Merge operation
def merge_judgments (J1 J2 : Judgment) : Judgment :=
  { coord := J1.coord
  , prop := J1.prop /\ J2.prop
  , carrier := J1.carrier \cup J2.carrier
  , evidence := J1.evidence |> J2.evidence
  , obligations := J1.obligations \cup J2.obligations
  , obstructions := J1.obstructions \cup J2.obstructions
  , trust := J1.trust |> J2.trust
  , provenance := J1.provenance ++ J2.provenance }

-- Merge commutativity and associativity
theorem merge_comm (J1 J2 : Judgment) :
    merge_judgments J1 J2 = merge_judgments J2 J1 := by
  simp [merge_judgments, And.comm, Set.union_comm,
        TrustAlgebra.join_comm]

theorem merge_assoc (J1 J2 J3 : Judgment) :
    merge_judgments (merge_judgments J1 J2) J3 =
    merge_judgments J1 (merge_judgments J2 J3) := by
  simp [merge_judgments, And.assoc, Set.union_assoc,
        TrustAlgebra.join_assoc]

-- Merge soundness
theorem merge_soundness
    (cover : JurisdictionCover S)
    (sheaves : forall d, JudgmentSheaf (cover.assign d))
    (compat : forall d1 d2,
      overlap_compatible (sheaves d1) (sheaves d2))
    : exists J : JudgmentSheaf S,
      H1_vanishes S J := by
  apply sheaf_gluing_from_cover cover sheaves compat
  exact cocycle_from_overlap_compat compat

-- Delegation chain bound
theorem delegation_chain_sound
    (chain : DelegationChain T)
    (h_len : chain.length <= 3)
    : chain.final_trust != T.contradicted := by
  induction chain.length with
  | zero => exact T.owner_not_contra
  | succ n ih =>
    apply T.attenuate_above_contra
    exact ih (Nat.le_of_succ_le h_len)

-- Jurisdiction lattice completeness
theorem jurisdiction_lattice_complete
    (S : Site) [Fintype S.Obj] (D : Finset Developer)
    : CompleteLattice (JurisdictionCover S D) :=
  complete_lattice_of_finite_join_and_meet
    (fun covers => coverage_union covers)
    (fun covers => coverage_meet_with_redistribution covers)

-- Cech cohomology vanishing implies merge
theorem cech_vanishing_implies_merge
    (cover : JurisdictionCover S)
    (h : H1_cech cover J = 0)
    : exists global : JudgmentSheaf S,
      forall d, restrict global (cover.assign d) =
        local_sheaf d := by
  exact sheaf_from_vanishing_H1 cover J h
\end{lstlisting}

% ------------------------------------------------------------------
\section{Implementation}\label{sec:implementation}
% ------------------------------------------------------------------

The team workflow system extends the \JuGeo{} Python implementation with three
new subsystems:

\paragraph{JurisdictionManager.}
Computes jurisdiction boundaries from code ownership metadata (CODEOWNERS files,
git blame). The site construction leverages \texttt{formal\_core\_site} with a
mean build time of \subFormalcoresiteSmallTime{} for small projects and
\subFormalcoresiteMediumTime{} for medium projects.

\paragraph{DelegationRegistry.}
Manages active delegation tokens with TTL enforcement. Built atop the trust
algebra, it uses \texttt{public\_alignment} to verify that delegations respect
the trust ordering. The alignment check runs in
\compProfPublicAlignmentMeanMs{}\,ms per call.

\paragraph{MergeEngine.}
Resolves concurrent judgment streams using Algorithm~\ref{alg:conflict}.
Conflict detection invokes \texttt{encode\_for\_solver} (mean time:
\compProfEncodeForSolverMeanMs{}\,ms) and dispatches to Z3 for satisfiability
checking.

The implementation adds approximately 1{,}200 lines of Python to the existing
codebase. The \texttt{discovery\_pipeline} subsystem
(\compProfDiscoveryPipelineMeanMs{}\,ms per invocation) identifies affected
coordinates, while \texttt{analogy\_transport}
(\compProfAnalogyTransportMeanMs{}\,ms) reuses judgments from structurally
similar scopes.

% ------------------------------------------------------------------
\section{Evaluation}\label{sec:evaluation}
% ------------------------------------------------------------------

\subsection{Experimental Setup}

We evaluate team workflows on \expTotalPrograms{} programs spanning
\expDomainCount{} domains. Each program is partitioned into 2--4 jurisdictions
simulating a multi-developer team. We measure merge overhead, conflict rate,
and end-to-end verification time.

All experiments use the \JuGeo{} Python implementation with the
\texttt{orchestrate\_verification} pipeline. The site is constructed via
\texttt{formal\_core\_site}, descent uses the eager strategy with depth limit~5,
and SMT queries are dispatched to Z3 via \texttt{encode\_for\_solver}. The
total experiment time across all benchmarks is \subTotalExperimentTime{}.

\subsection{Per-Domain Results}

\begin{table}[t]
\centering
\caption{Team workflow verification results across domains.}
\label{tab:team-results}
\begin{tabular}{lrrrrr}
\toprule
\textbf{Domain} & \textbf{Programs} & \textbf{Verified} &
  \textbf{Avg Coords} & \textbf{Avg Props} & \textbf{Conflicts} \\
\midrule
Data Structures & \expDomainDsCount{} & \expDomainDsVerified{} &
  \expDomainDsAvgCoords{} & \expDomainDsAvgProps{} & 0 \\
Math            & \expDomainMathCount{} & \expDomainMathVerified{} &
  \expDomainMathAvgCoords{} & \expDomainMathAvgProps{} & 0 \\
State Machines  & \expDomainSmCount{} & \expDomainSmVerified{} &
  \expDomainSmAvgCoords{} & \expDomainSmAvgProps{} & 0 \\
Sorting         & \expDomainSortCount{} & \expDomainSortVerified{} &
  \expDomainSortAvgCoords{} & \expDomainSortAvgProps{} & 0 \\
String          & \expDomainStrCount{} & \expDomainStrVerified{} &
  \expDomainStrAvgCoords{} & \expDomainStrAvgProps{} & 0 \\
Web             & \expDomainWebCount{} & \expDomainWebVerified{} &
  \expDomainWebAvgCoords{} & \expDomainWebAvgProps{} & 0 \\
\midrule
\textbf{Total}  & \expTotalPrograms{} & \expVerified{} &
  \expCoordsMean{} & \expPropsMean{} & \expObstructionTotal{} \\
\bottomrule
\end{tabular}
\end{table}

Table~\ref{tab:team-results} confirms that all \expTotalPrograms{} programs
verify with zero merge conflicts. The overall accuracy is \expAccuracy{} with
\expPropsSum{} total propositions (\expPropsOkSum{} passing).

\subsection{Per-Domain Timing Analysis}

\begin{table}[t]
\centering
\caption{Per-domain timing statistics for team verification.}
\label{tab:domain-timing}
\begin{tabular}{lrrr}
\toprule
\textbf{Domain} & \textbf{Avg Time} & \textbf{Min Time} & \textbf{Max Time} \\
\midrule
Sorting         & \suppDomainSortAvgTime{} & \suppDomainSortMinTime{}
                & \suppDomainSortMaxTime{} \\
Data Structures & \suppDomainDsAvgTime{} & \suppDomainDsMinTime{}
                & \suppDomainDsMaxTime{} \\
Math            & \suppDomainMathAvgTime{} & \suppDomainMathMinTime{}
                & \suppDomainMathMaxTime{} \\
String          & \suppDomainStrAvgTime{} & \suppDomainStrMinTime{}
                & \suppDomainStrMaxTime{} \\
Web             & \suppDomainWebAvgTime{} & \suppDomainWebMinTime{}
                & \suppDomainWebMaxTime{} \\
State Machines  & \suppDomainSmAvgTime{} & \suppDomainSmMinTime{}
                & \suppDomainSmMaxTime{} \\
\midrule
\textbf{Overall}& \expTimeMean{} & \expTimeMin{} & \expTimeMax{} \\
\bottomrule
\end{tabular}
\end{table}

Table~\ref{tab:domain-timing} breaks down timing by domain.
State Machines exhibit the highest average (\suppDomainSmAvgTime{}) due to
larger state spaces; Strings are fastest (\suppDomainStrAvgTime{}). The overall
mean of \expTimeMean{} (median \expTimeMedian{}) fits interactive feedback.

\subsection{Proposition Breakdown}

\begin{table}[t]
\centering
\caption{Proposition kind distribution across all verified programs.}
\label{tab:prop-breakdown}
\begin{tabular}{lrr}
\toprule
\textbf{Proposition Kind} & \textbf{Count} & \textbf{Percentage} \\
\midrule
Structural  & \suppPropStructuralCount{} & \suppPropStructuralPct{} \\
Behavioral  & \suppPropBehavioralCount{} & \suppPropBehavioralPct{} \\
Relational  & \suppPropRelationalCount{} & \suppPropRelationalPct{} \\
Resource    & \suppPropResourceCount{} & \suppPropResourcePct{} \\
\midrule
\textbf{Total} & \suppPropTotal{} & 100\% \\
\bottomrule
\end{tabular}
\end{table}

Behavioral propositions dominate (\suppPropBehavioralPct{}), reflecting
functional correctness emphasis. Structural propositions
(\suppPropStructuralPct{}) capture type invariants, while relational
(\suppPropRelationalPct{}) and resource (\suppPropResourcePct{}) propositions
are less frequent.

\subsection{Coordinate Distribution}

\begin{table}[t]
\centering
\caption{Coordinate kind distribution by granularity level.}
\label{tab:coord-breakdown}
\begin{tabular}{lrr}
\toprule
\textbf{Coordinate Kind} & \textbf{Count} & \textbf{Percentage} \\
\midrule
Module    & \suppCoordModuleCount{} & \suppCoordModulePct{} \\
Function  & \suppCoordFunctionCount{} & \suppCoordFunctionPct{} \\
Interface & \suppCoordInterfaceCount{} & \suppCoordInterfacePct{} \\
\midrule
\textbf{Total} & \suppCoordTotal{} & 100\% \\
\bottomrule
\end{tabular}
\end{table}

Function-level coordinates dominate (\suppCoordFunctionPct{}),
indicating that most verification targets individual functions. Module
coordinates (\suppCoordModulePct{}) and interface coordinates
(\suppCoordInterfacePct{}) are less frequent. This distribution informs
jurisdiction granularity: function-level jurisdictions provide the finest
parallelism.

\subsection{Size Scaling}

\begin{table}[t]
\centering
\caption{Team verification scaling with program size.}
\label{tab:size-scaling}
\begin{tabular}{lrrrrr}
\toprule
\textbf{Size} & \textbf{Count} & \textbf{Lines} &
  \textbf{Coords} & \textbf{Props} & \textbf{Time} \\
\midrule
Tiny   & \compTinyCount{} & \compTinyMeanLines{}
       & \compTinyMeanCoords{} & \compTinyMeanProps{} & \compTinyMeanTime{} \\
Small  & \compSmallCount{} & \compSmallMeanLines{}
       & \compSmallMeanCoords{} & \compSmallMeanProps{} & \compSmallMeanTime{} \\
Medium & \compMediumCount{} & \compMediumMeanLines{}
       & \compMediumMeanCoords{} & \compMediumMeanProps{} & \compMediumMeanTime{} \\
Large  & \compLargeCount{} & \compLargeMeanLines{}
       & \compLargeMeanCoords{} & \compLargeMeanProps{} & \compLargeMeanTime{} \\
\bottomrule
\end{tabular}
\end{table}

Notably, verification time does not increase monotonically with program size:
large programs (\compLargeMeanTime{}) verify faster than tiny programs
(\compTinyMeanTime{}) because larger programs have more exploitable structure.

\subsection{Developer Scaling}

\begin{table}[t]
\centering
\caption{Site construction time scaling with number of coordinates.}
\label{tab:dev-scaling}
\begin{tabular}{lr}
\toprule
\textbf{Coordinates} & \textbf{Construction Time} \\
\midrule
2  & \subScaleTwoTime{} \\
4  & \subScaleFourTime{} \\
8  & \subScaleEightTime{} \\
12 & \subScaleTwelveTime{} \\
16 & \subScaleSixteenTime{} \\
20 & \subScaleTwentyTime{} \\
\bottomrule
\end{tabular}
\end{table}

Table~\ref{tab:dev-scaling} shows near-constant site construction time as
coordinates grow from 2 to 20, confirming that jurisdiction partitioning
(Algorithm~\ref{alg:jurisdiction}) introduces negligible overhead.

\subsection{Descent Strategy Comparison}

\begin{table}[t]
\centering
\caption{Descent strategy comparison for team verification.}
\label{tab:descent-comparison}
\begin{tabular}{lrrr}
\toprule
\textbf{Strategy} & \textbf{Runs} & \textbf{Success Rate} & \textbf{Time} \\
\midrule
Eager       & \compDescentEagerRuns{} & \compDescentEagerRate{}
            & \subDescentEagerTime{} \\
Exhaustive  & \compDescentExhaustiveRuns{} & \compDescentExhaustiveRate{}
            & \subDescentExhaustiveTime{} \\
Iterative   & \compDescentIterativeRuns{} & \compDescentIterativeRate{}
            & \subDescentIterativeTime{} \\
Optimistic  & \compDescentOptimisticRuns{} & \compDescentOptimisticRate{}
            & \subDescentOptimisticTime{} \\
\bottomrule
\end{tabular}
\end{table}

All four descent strategies achieve 100\% success rate. We default to eager
(\subDescentEagerTime{}) for team workflows, with exhaustive
(\subDescentExhaustiveTime{}) as a fallback for overlap zones.

\subsection{Cache Effectiveness}

The team workflow benefits significantly from judgment caching.
When a developer re-verifies after a small change, cached judgments for
unchanged scopes are reused.

\begin{table}[t]
\centering
\caption{Cache performance for incremental team verification.}
\label{tab:cache}
\begin{tabular}{lr}
\toprule
\textbf{Metric} & \textbf{Value} \\
\midrule
Cold start mean time  & \suppColdMeanMs{}\,ms \\
Warm cache mean time  & \suppWarmMeanMs{}\,ms \\
Cache speedup         & \suppCacheSpeedup{} \\
\bottomrule
\end{tabular}
\end{table}

The \suppCacheSpeedup{} speedup (Table~\ref{tab:cache}) makes incremental
PR review workflows practical: verification completes in \suppWarmMeanMs{}\,ms
on warm cache, well before a reviewer finishes reading the diff.

\subsection{Morphism and Trust Analysis}

\paragraph{Merge Overhead.}
The merge operation adds negligible overhead: with \compMorphTotal{} total
morphisms (\compMorphRestriction{} restrictions at \compMorphRestrictionPct{},
\compMorphTransport{} transports at \compMorphTransportPct{}, and
\compMorphInclusion{} inclusions at \compMorphInclusionPct{}), the merge
processes all overlap morphisms in under \compTimeMean{} per program.

\paragraph{Trust Distribution.}
All \compTrustTotal{} trust annotations resolve to $\Tsolver$-discharged
(\compTrustSolverdischargedPct{}), with \compTrustUnverified{} unverified
(\compTrustUnverifiedPct{}). This confirms that team delegation does not
degrade verification quality.

\paragraph{Equivalence Checking.}
All \compEquivBothVerified{} of \compEquivPairs{} equivalence pairs verify
consistently (\compEquivSameStructure{} same structure, mean time
\compEquivMeanTime{}), validating that equivalent jurisdictions produce
compatible judgments.

\subsection{Refinement Classification}

\begin{table}[t]
\centering
\caption{Refinement pair classification for jurisdiction comparisons.}
\label{tab:refinement}
\begin{tabular}{lrr}
\toprule
\textbf{Relation} & \textbf{Pairs} & \textbf{Both OK} \\
\midrule
Forward refinement   & \suppForwardPairs{} & \suppForwardBothOk{} \\
Equivalence          & \suppEquivPairs{} & \suppEquivBothOk{} \\
Incomparable         & \suppIncompPairs{} & \suppIncompBothOk{} \\
\midrule
\textbf{Total}       & \suppTotalPairs{} & \suppTotalPairs{} \\
\bottomrule
\end{tabular}
\end{table}

Table~\ref{tab:refinement} confirms all refinement pairs verify
regardless of ordering, showing the jurisdiction lattice is compatible with
verification: refining to a smaller scope preserves correctness.

\subsection{Bug Detection Under Team Workflows}

\paragraph{Bug Detection.}
On \compBuggyCount{} deliberately buggy variants, the team workflow detects all
bugs (\compBuggyFullPass{} full detections, mean time \compBuggyMeanTime{},
pass ratio \compBuggyMeanPropRatio{}). Even bugs within a single jurisdiction
correctly propagate as obstructions at the global level.

\subsection{Subsystem Profiling}

\begin{table}[t]
\centering
\caption{Key subsystem profiling for team workflow operations.}
\label{tab:subsystem-profile}
\begin{tabular}{lrr}
\toprule
\textbf{Subsystem} & \textbf{Mean Time (ms)} & \textbf{Success Rate} \\
\midrule
\texttt{orchestrate\_verification} & \compProfOrchestrateVerificationMeanMs{}
  & \compProfOrchestrateVerificationRate{} \\
\texttt{encode\_for\_solver} & \compProfEncodeForSolverMeanMs{}
  & \compProfEncodeForSolverRate{} \\
\texttt{discovery\_pipeline} & \compProfDiscoveryPipelineMeanMs{}
  & \compProfDiscoveryPipelineRate{} \\
\texttt{public\_alignment} & \compProfPublicAlignmentMeanMs{}
  & \compProfPublicAlignmentRate{} \\
\texttt{formal\_core\_site} & \compProfFormalCoreSiteMeanMs{}
  & \compProfFormalCoreSiteRate{} \\
\texttt{analogy\_transport} & \compProfAnalogyTransportMeanMs{}
  & \compProfAnalogyTransportRate{} \\
\texttt{run\_full\_descent} & \compProfRunFullDescentMeanMs{}
  & \compProfRunFullDescentRate{} \\
\texttt{judgment\_sheaf} & \compProfJudgmentSheafMeanMs{}
  & \compProfJudgmentSheafRate{} \\
\texttt{trust\_presheaf} & \compProfTrustPresheafMeanMs{}
  & \compProfTrustPresheafRate{} \\
\texttt{repair\_semantics} & \compProfRepairSemanticsMeanMs{}
  & \compProfRepairSemanticsRate{} \\
\bottomrule
\end{tabular}
\end{table}

Table~\ref{tab:subsystem-profile} profiles key subsystems. Orchestration
(\compProfOrchestrateVerificationMeanMs{}\,ms) dominates, while SMT encoding
(\compProfEncodeForSolverMeanMs{}\,ms) and site construction
(\compProfFormalCoreSiteMeanMs{}\,ms) are sub-millisecond.

% ------------------------------------------------------------------
\section{Related Work}\label{sec:related}
% ------------------------------------------------------------------

\paragraph{Collaborative Verification.}
VeriFast~\cite{jacobs2011verifast} and Viper~\cite{muller2016viper} support
modular verification but lack explicit team coordination. Dafny's
module system~\cite{leino2010dafny} provides some ownership semantics but
without trust-attenuated delegation. Iris~\cite{jung2018iris} introduces
higher-order concurrent separation logic with ghost state, supporting modular
reasoning about shared resources, but does not address team-level jurisdiction
or authority delegation. Our framework lifts these modular verification ideas
to the organizational level, where the ``modules'' are developer jurisdictions
and the ``interfaces'' are overlap zones with cohomological compatibility
conditions.

\paragraph{Code Ownership Models.}
GitHub's CODEOWNERS and Google's code review policies define ownership at the
file level. Our jurisdiction model lifts this to the semantic site level,
enabling finer-grained, scope-aware ownership. Bird et al.~\cite{bird2011ownership}
study the relationship between code ownership and software quality, finding
that diffuse ownership correlates with defects. Our jurisdiction lattice
(Definition~\ref{thm:juris-lattice}) formalizes the spectrum from minimal
partitions (single owner) to maximal overlap (shared ownership), enabling
teams to choose the granularity that balances parallelism against coordination
cost.

\paragraph{Merge Verification.}
Semantic merge tools like SemanticMerge~\cite{semanticmerge} resolve textual
conflicts but do not verify correctness. Our approach merges \emph{judgments},
ensuring that merged code retains its verification status. Sousa et
al.~\cite{sousa2018incremental} study incremental verification after code
changes, maintaining verification results across edits. Our merge engine
extends this to concurrent edits by multiple developers, using sheaf-theoretic
gluing rather than incremental re-checking.

\paragraph{Sheaf-Theoretic Approaches.}
Goguen~\cite{goguen1992sheaves} used sheaves for modeling information systems.
Robinson~\cite{robinson2014topological} applies sheaf theory to data fusion.
Curry~\cite{curry2014sheaves} uses sheaves for information flow in networks.
Our work applies these ideas to verification, where ``local data'' are judgments
and ``consistency'' is overlap compatibility.

\paragraph{\v{C}ech Cohomology in Software.}
Cohomological methods are established in algebraic
topology~\cite{hartshorne1977algebraic,bott1982differential} but novel in
software engineering. Goubault~\cite{goubault2000geometry} and
Fajstrup et al.~\cite{fajstrup2016directed} use directed algebraic topology
for concurrency, where deadlocks correspond to nontrivial homology. We use
\v{C}ech cohomology specifically for merge analysis.

\paragraph{Trust and Authorization.}
SPKI/SDSI~\cite{ellison1999spki} and capability systems~\cite{miller2006robust}
model authorization with delegation chains. Our $\DelegCat$ operates in $\Talg$
with strict attenuation, analogous to capability attenuation.

\paragraph{Distributed Verification.}
Flyspeck~\cite{hales2017formal} and seL4~\cite{klein2009sel4} involve large
teams but coordinate informally. Our framework provides formal consistency
guarantees via the sheaf condition.

% ------------------------------------------------------------------
\section{Conclusion}\label{sec:conclusion}
% ------------------------------------------------------------------

We have presented a framework for multi-developer workflows in \JuGeo{},
grounded in jurisdiction partitioning and trust-attenuated authority delegation.
The merge semantics for concurrent judgment streams are sound by construction,
leveraging the sheaf-theoretic foundations of the judgment-geometric approach.
Our evaluation on \expTotalPrograms{} programs confirms zero merge-induced
obstructions and negligible overhead.

The key theoretical contribution is the identification of team merge soundness
with the vanishing of the first \v{C}ech cohomology group
$\Hcoh{1}(\mathcal{U}, \J) = 0$ over the jurisdiction cover. This
cohomological criterion provides a clean, compositional test: individual
developers verify their jurisdictions independently, the merge engine checks
pairwise compatibility on overlaps, and the vanishing of $\Hcoh{1}$ guarantees
global consistency. The spectral sequence for hierarchical teams enables
multi-stage verification in large organizations.

From a practical standpoint, the framework introduces minimal overhead. The
jurisdiction partitioning algorithm runs in sub-millisecond time, the merge
engine leverages the existing descent and SMT infrastructure, and the
\suppCacheSpeedup{} cache speedup makes incremental team verification
interactive. The delegation protocol, with its bounded chain length and
automatic revocation propagation, provides a lightweight authorization layer
that integrates naturally with existing code ownership tools.

Future work includes:
\begin{enumerate}[leftmargin=*]
  \item \textbf{Distributed teams with network partitions.} Extending the
        merge semantics to handle delayed synchronization, where developers
        may produce conflicting judgments that are only detected at merge time.
        The higher cohomology groups $\Hcoh{p}$ for $p \geq 2$ may play a
        role in classifying higher-order inconsistencies.
  \item \textbf{Git integration.} Implementing jurisdiction-aware merge
        drivers that invoke \JuGeo{} verification as part of \texttt{git merge},
        providing semantic merge conflict detection alongside textual merge.
  \item \textbf{Continuous verification.} Running the team workflow
        continuously in CI/CD pipelines, maintaining a live judgment sheaf
        that updates incrementally as developers push changes.
  \item \textbf{Formal verification of the framework itself.} Completing the
        \LEAN{} formalization of all theorems presented in this paper, producing
        a fully machine-checked proof of team merge soundness.
  \item \textbf{Higher-order team structures.} Modeling nested teams
        (teams of teams) using higher categorical constructs---specifically,
        2-sheaves on a 2-site of organizational hierarchies.
\end{enumerate}

\bibliographystyle{alpha}
\begin{thebibliography}{99}
\bibitem{jacobs2011verifast}
B.~Jacobs, J.~Smans, P.~Philippaerts, F.~Piessens, W.~Penninckx, and
F.~Vogels.
\newblock VeriFast: A powerful, sound, predictable, fast verifier for C and
Java.
\newblock In \emph{NFM}, 2011.

\bibitem{muller2016viper}
P.~M{\"u}ller, M.~Schwerhoff, and A.~J.~Summers.
\newblock Viper: A verification infrastructure for permission-based reasoning.
\newblock In \emph{VMCAI}, 2016.

\bibitem{leino2010dafny}
K.~R.~M.~Leino.
\newblock Dafny: An automatic program verifier for functional correctness.
\newblock In \emph{LPAR}, 2010.

\bibitem{semanticmerge}
Codice Software.
\newblock Semantic{M}erge: The language-aware merge tool, 2020.

\bibitem{jugeo-paper19}
\JuGeo{} Development Team.
\newblock Authority delegation in judgment-geometric verification.
\newblock Paper~19 in the \JuGeo{} series.

\bibitem{jugeo-paper27}
\JuGeo{} Development Team.
\newblock Certificate chains for compositional verification.
\newblock Paper~27 in the \JuGeo{} series.

\bibitem{jugeo-paper37}
\JuGeo{} Development Team.
\newblock Pack federation and distributed judgment synthesis.
\newblock Paper~37 in the \JuGeo{} series.

\bibitem{henzinger1998you}
T.~A.~Henzinger, S.~Qadeer, and S.~K.~Rajamani.
\newblock You assume, we guarantee: Methodology and case studies.
\newblock In \emph{CAV}, 1998.

\bibitem{clarke2000counterexample}
E.~Clarke, O.~Grumberg, S.~Jha, Y.~Lu, and H.~Veith.
\newblock Counterexample-guided abstraction refinement.
\newblock In \emph{CAV}, 2000.

\bibitem{jung2018iris}
R.~Jung, R.~Krebbers, J.-H.~Jourdan, A.~Bizjak, L.~Birkedal, and
D.~Dreyer.
\newblock Iris from the ground up: A modular foundation for higher-order
concurrent separation logic.
\newblock \emph{J.~Funct.~Program.}, 28:e20, 2018.

\bibitem{bird2011ownership}
C.~Bird, N.~Nagappan, B.~Murphy, H.~Gall, and P.~Devanbu.
\newblock Don't touch my code! {E}xamining the effects of ownership on
software quality.
\newblock In \emph{ESEC/FSE}, 2011.

\bibitem{sousa2018incremental}
M.~Sousa and I.~Dillig.
\newblock Cartesian {H}oare logic for verifying $k$-safety properties.
\newblock In \emph{PLDI}, 2016.

\bibitem{goguen1992sheaves}
J.~A.~Goguen.
\newblock Sheaf semantics for concurrent interacting objects.
\newblock \emph{Math.~Struct.~Comput.~Sci.}, 2(2):159--191, 1992.

\bibitem{robinson2014topological}
M.~Robinson.
\newblock \emph{Topological Signal Processing}.
\newblock Springer, 2014.

\bibitem{curry2014sheaves}
J.~M.~Curry.
\newblock Sheaves, cosheaves and applications.
\newblock Ph.D. thesis, University of Pennsylvania, 2014.

\bibitem{hartshorne1977algebraic}
R.~Hartshorne.
\newblock \emph{Algebraic Geometry}.
\newblock Springer, 1977.

\bibitem{bott1982differential}
R.~Bott and L.~W.~Tu.
\newblock \emph{Differential Forms in Algebraic Topology}.
\newblock Springer, 1982.

\bibitem{goubault2000geometry}
{\'E}.~Goubault.
\newblock Geometry and concurrency: A user's guide.
\newblock \emph{Math.~Struct.~Comput.~Sci.}, 10(4):411--425, 2000.

\bibitem{fajstrup2016directed}
L.~Fajstrup, E.~Goubault, E.~Haucourt, S.~Mimram, and M.~Raussen.
\newblock \emph{Directed Algebraic Topology and Concurrency}.
\newblock Springer, 2016.

\bibitem{ellison1999spki}
C.~Ellison, B.~Frantz, B.~Lampson, R.~Rivest, B.~Thomas, and T.~Ylonen.
\newblock {SPKI} certificate theory.
\newblock RFC 2693, IETF, 1999.

\bibitem{miller2006robust}
M.~S.~Miller, K.-P.~Yee, and J.~Shapiro.
\newblock Capability myths demolished.
\newblock Technical Report SRL2003-02, Johns Hopkins University, 2003.

\bibitem{hales2017formal}
T.~C.~Hales et~al.
\newblock A formal proof of the {K}epler conjecture.
\newblock \emph{Forum Math. Pi}, 5:e2, 2017.

\bibitem{klein2009sel4}
G.~Klein et~al.
\newblock {seL4}: Formal verification of an {OS} kernel.
\newblock In \emph{SOSP}, 2009.
\end{thebibliography}

\end{document}
